% ============================================================
%  SAMURAI · МПС, объяснённая с нуля
%
%  Курс-учебник по модели пространства состояний (МПС) робота
%  SAMURAI. Story-driven: от «нажал кнопку — поехал» к матрицам
%  A, B и QP MPC-регулятора.
%
%  Компиляция: pdflatex main.tex  (два прохода для оглавления)
%  Каждая глава из chapters/ может собираться отдельно — см.
%  README.md в этой папке.
% ============================================================
\documentclass[12pt, a4paper, openany]{book}

% ============================================================
%  SAMURAI · МПС, объяснённая с нуля
%  Общий preamble (пакеты, окружения, макросы).
%
%  Файл подключается из main.tex и из «mini-main» каждой главы,
%  чтобы любую главу можно было собрать в отдельный PDF.
% ============================================================

% ── Кодировка и язык ─────────────────────────────────────────────
\usepackage[utf8]{inputenc}
\usepackage[T2A]{fontenc}
\usepackage[russian, english]{babel}

% ── Геометрия страницы ───────────────────────────────────────────
\usepackage[
  top=25mm, bottom=25mm,
  left=30mm, right=20mm
]{geometry}

% ── Математика ───────────────────────────────────────────────────
\usepackage{amsmath, amssymb, amsthm}
\usepackage{mathtools}
\usepackage{bm}

% ── Графика и цвет ───────────────────────────────────────────────
\usepackage{xcolor}
\usepackage{graphicx}
\usepackage{tikz}
\usetikzlibrary{arrows.meta, positioning, shapes.geometric, calc, fit,
                decorations.pathmorphing, decorations.markings, backgrounds}

% ── Код ──────────────────────────────────────────────────────────
\usepackage{listings}
\usepackage{courier}

\definecolor{codebg}{RGB}{245,247,250}
\definecolor{codeframe}{RGB}{180,190,210}
\definecolor{keyword}{RGB}{0,100,200}
\definecolor{comment}{RGB}{100,130,100}
\definecolor{string}{RGB}{163,21,21}
\definecolor{number}{RGB}{9,134,88}

% Поддержка кириллицы в listings.
\lstset{
  extendedchars=true,
  inputencoding=utf8,
  literate=
    {а}{{\selectfont\char224}}1{б}{{\selectfont\char225}}1
    {в}{{\selectfont\char226}}1{г}{{\selectfont\char227}}1
    {д}{{\selectfont\char228}}1{е}{{\selectfont\char229}}1
    {ё}{{\selectfont\char184}}1{ж}{{\selectfont\char230}}1
    {з}{{\selectfont\char231}}1{и}{{\selectfont\char232}}1
    {й}{{\selectfont\char233}}1{к}{{\selectfont\char234}}1
    {л}{{\selectfont\char235}}1{м}{{\selectfont\char236}}1
    {н}{{\selectfont\char237}}1{о}{{\selectfont\char238}}1
    {п}{{\selectfont\char239}}1{р}{{\selectfont\char240}}1
    {с}{{\selectfont\char241}}1{т}{{\selectfont\char242}}1
    {у}{{\selectfont\char243}}1{ф}{{\selectfont\char244}}1
    {х}{{\selectfont\char245}}1{ц}{{\selectfont\char246}}1
    {ч}{{\selectfont\char247}}1{ш}{{\selectfont\char248}}1
    {щ}{{\selectfont\char249}}1{ъ}{{\selectfont\char250}}1
    {ы}{{\selectfont\char251}}1{ь}{{\selectfont\char252}}1
    {э}{{\selectfont\char253}}1{ю}{{\selectfont\char254}}1
    {я}{{\selectfont\char255}}1
    {А}{{\selectfont\char192}}1{Б}{{\selectfont\char193}}1
    {В}{{\selectfont\char194}}1{Г}{{\selectfont\char195}}1
    {Д}{{\selectfont\char196}}1{Е}{{\selectfont\char197}}1
    {Ё}{{\selectfont\char168}}1{Ж}{{\selectfont\char198}}1
    {З}{{\selectfont\char199}}1{И}{{\selectfont\char200}}1
    {Й}{{\selectfont\char201}}1{К}{{\selectfont\char202}}1
    {Л}{{\selectfont\char203}}1{М}{{\selectfont\char204}}1
    {Н}{{\selectfont\char205}}1{О}{{\selectfont\char206}}1
    {П}{{\selectfont\char207}}1{Р}{{\selectfont\char208}}1
    {С}{{\selectfont\char209}}1{Т}{{\selectfont\char210}}1
    {У}{{\selectfont\char211}}1{Ф}{{\selectfont\char212}}1
    {Х}{{\selectfont\char213}}1{Ц}{{\selectfont\char214}}1
    {Ч}{{\selectfont\char215}}1{Ш}{{\selectfont\char216}}1
    {Щ}{{\selectfont\char217}}1{Ъ}{{\selectfont\char218}}1
    {Ы}{{\selectfont\char219}}1{Ь}{{\selectfont\char220}}1
    {Э}{{\selectfont\char221}}1{Ю}{{\selectfont\char222}}1
    {Я}{{\selectfont\char223}}1
    {—}{{\textemdash}}1{–}{{\textendash}}1{«}{{\guillemotleft}}1
    {»}{{\guillemotright}}1{№}{{N\textsuperscript{o}}}1
    {²}{{\textsuperscript{2}}}1{·}{{$\cdot$}}1{°}{{$^\circ$}}1
    {≤}{{$\leq$}}1{≥}{{$\geq$}}1{≠}{{$\neq$}}1{≈}{{$\approx$}}1
    {±}{{$\pm$}}1{×}{{$\times$}}1{÷}{{$\div$}}1{√}{{$\surd$}}1
    {→}{{$\to$}}1{←}{{$\leftarrow$}}1{⇒}{{$\Rightarrow$}}1
    {⇔}{{$\Leftrightarrow$}}1
    {θ}{{$\theta$}}1{φ}{{$\varphi$}}1{ω}{{$\omega$}}1
    {Φ}{{$\Phi$}}1{Γ}{{$\Gamma$}}1{Σ}{{$\Sigma$}}1{Ω}{{$\Omega$}}1
    {ε}{{$\varepsilon$}}1{σ}{{$\sigma$}}1{μ}{{$\mu$}}1{ν}{{$\nu$}}1
    {ρ}{{$\rho$}}1{δ}{{$\delta$}}1{η}{{$\eta$}}1{ψ}{{$\psi$}}1
    {α}{{$\alpha$}}1{β}{{$\beta$}}1{γ}{{$\gamma$}}1
    {τ}{{$\tau$}}1{π}{{$\pi$}}1{Δ}{{$\Delta$}}1{λ}{{$\lambda$}}1
}

\lstdefinestyle{pythonstyle}{
  language=Python,
  backgroundcolor=\color{codebg},
  frame=single, rulecolor=\color{codeframe}, framesep=4pt,
  basicstyle=\footnotesize\ttfamily,
  keywordstyle=\color{keyword}\bfseries,
  commentstyle=\color{comment}\itshape,
  stringstyle=\color{string},
  numberstyle=\scriptsize\ttfamily\color{gray},
  numbers=left, numbersep=10pt, xleftmargin=2.5em,
  showstringspaces=false, breaklines=true,
  tabsize=4, captionpos=b,
  morekeywords={True, False, None, self},
}

\lstdefinestyle{matlabstyle}{
  language=Matlab,
  backgroundcolor=\color{codebg},
  frame=single, rulecolor=\color{codeframe}, framesep=4pt,
  basicstyle=\footnotesize\ttfamily,
  keywordstyle=\color{keyword}\bfseries,
  commentstyle=\color{comment}\itshape,
  stringstyle=\color{string},
  numberstyle=\scriptsize\ttfamily\color{gray},
  numbers=left, numbersep=10pt, xleftmargin=2.5em,
  showstringspaces=false, breaklines=true,
  tabsize=2, captionpos=b,
}

\lstdefinestyle{tsstyle}{
  language=Java,                % близкий к TS лексер
  backgroundcolor=\color{codebg},
  frame=single, rulecolor=\color{codeframe}, framesep=4pt,
  basicstyle=\footnotesize\ttfamily,
  keywordstyle=\color{keyword}\bfseries,
  commentstyle=\color{comment}\itshape,
  stringstyle=\color{string},
  numberstyle=\scriptsize\ttfamily\color{gray},
  numbers=left, numbersep=10pt, xleftmargin=2.5em,
  showstringspaces=false, breaklines=true,
  tabsize=2, captionpos=b,
  morekeywords={const, let, function, export, import, from, type,
                interface, number, string, boolean, return, if, else},
}

\lstset{style=pythonstyle}

% ── Таблицы ──────────────────────────────────────────────────────
\usepackage{booktabs}
\usepackage{array}
\usepackage{longtable}
\usepackage{float}

% ── Гиперссылки ──────────────────────────────────────────────────
\usepackage[
  colorlinks=true,
  linkcolor=blue!70!black,
  citecolor=green!60!black,
  urlcolor=blue!80!black,
  bookmarks=true,
  bookmarksnumbered=true
]{hyperref}

% ── Колонтитулы ──────────────────────────────────────────────────
\usepackage{fancyhdr}
\setlength{\headheight}{14pt}     % fancyhdr требует >= 13.6pt для русского fontenc
\addtolength{\topmargin}{-2pt}
\pagestyle{fancy}
\fancyhf{}
\fancyhead[LE,RO]{\small\thepage}
\fancyhead[LO]{\small\nouppercase{\rightmark}}
\fancyhead[RE]{\small\nouppercase{\leftmark}}
\renewcommand{\headrulewidth}{0.4pt}

% ── Параграфы ────────────────────────────────────────────────────
\usepackage{parskip}
\setlength{\parindent}{0pt}

% ── Окружения ────────────────────────────────────────────────────
\usepackage{tcolorbox}
\tcbuselibrary{skins, breakable}

% «Идея»: интуиция/мотивация перед формальной частью.
\newtcolorbox{ideabox}[1][]{
  enhanced, breakable,
  colback=blue!4!white, colframe=blue!55!black,
  fonttitle=\bfseries, title=Идея, #1
}
% «Формула»: ключевое уравнение, на которое потом ссылаемся.
\newtcolorbox{formulabox}[1][]{
  enhanced, breakable,
  colback=blue!3!white, colframe=blue!50!black,
  fonttitle=\bfseries, title=Формула, #1
}
% «Примечание»: оранжевая «полочка», объяснение в сторону.
\newtcolorbox{notebox}[1][]{
  enhanced,
  colback=orange!8!white, colframe=orange!70!black,
  fonttitle=\bfseries, title=Примечание, #1
}
% «Внимание»: красная коробка — типичная ошибка / подводный камень.
\newtcolorbox{warnbox}[1][]{
  enhanced,
  colback=red!4!white, colframe=red!70!black,
  fonttitle=\bfseries, title=Внимание, #1
}
% «В коде»: серая «полочка» — ссылка на файл/строку с пояснением.
\newtcolorbox{codebox}[1][]{
  enhanced, breakable,
  colback=gray!6!white, colframe=gray!50!black,
  fonttitle=\bfseries\ttfamily, title=В коде, #1
}
% «Пример»: численный пример (зелёная).
\newtcolorbox{examplebox}[1][]{
  enhanced, breakable,
  colback=green!4!white, colframe=green!45!black,
  fonttitle=\bfseries, title=Пример, #1
}
% «Что если?»: фиолетовая — учебное «а что будет, если...».
\newtcolorbox{whatifbox}[1][]{
  enhanced, breakable,
  colback=violet!5!white, colframe=violet!55!black,
  fonttitle=\bfseries, title=Что если?, #1
}

% ── Нумерация уравнений по главам ────────────────────────────────
\numberwithin{equation}{chapter}

% ── Макросы ──────────────────────────────────────────────────────
\newcommand{\filepath}[1]{\texttt{\small\color{gray!80!black}#1}}
\newcommand{\vect}[1]{\bm{#1}}
\newcommand{\mat}[1]{\mathbf{#1}}
\newcommand{\T}{^{\mathsf{T}}}
\newcommand{\Real}{\mathbb{R}}
\newcommand{\diag}{\operatorname{diag}}
\newcommand{\rank}{\operatorname{rank}}
\newcommand{\sign}{\operatorname{sign}}
\DeclareMathOperator*{\argmin}{arg\,min}
\DeclareMathOperator*{\argmax}{arg\,max}
% \half — короткое 1/2.
\newcommand{\half}{\tfrac{1}{2}}


% ============================================================
\begin{document}

% ── Титульная страница ───────────────────────────────────────────
\begin{titlepage}
\centering
\vspace*{2cm}
{\Huge\bfseries SAMURAI Robot}\\[0.4cm]
{\Large\bfseries МПС, объяснённая с нуля}\\[0.3cm]
{\large Модель пространства состояний и MPC-регулятор\\
        в коде, на пальцах и в формулах}\\[1cm]
\rule{\textwidth}{1pt}\\[0.5cm]
{\large Учебник-сопровождение к ветке \texttt{feat/mps}}\\[0.5cm]
\rule{\textwidth}{1pt}\\[2cm]

\begin{tabular}{ll}
\textbf{Целевая аудитория:} & студент 3--4 курса ИТ-вуза\\
\textbf{Что нужно знать:}   & матан, линейная алгебра, основы Python\\
\textbf{Что объясняется:}   & ТАУ, дифференциальный привод, ЛКР, MPC\\
\textbf{Робот:}             & гусеничное шасси, 2 мотора, Raspberry Pi 4\\
\textbf{Размерность модели:}& $n = 5$ (состояние), $r = 2$ (управление)\\
\textbf{Период дискретизации:} & $T_s = 0.02$ с (50 Гц, секция \texttt{mps:})\\
\textbf{Стиль изложения:}   & story-driven --- от сценария к формуле\\
\end{tabular}

\vfill
{\large Основано на:\\
Козлов В.\,Н., Куприянов В.\,Е., Шашихин В.\,Н.\\
\textit{«Теория автоматического управления»} (СПбГПУ, 2008)}\\[0.4cm]
{\large \today}
\end{titlepage}

% ── Оглавление ───────────────────────────────────────────────────
\tableofcontents
\newpage

% ── Главы ────────────────────────────────────────────────────────
% ============================================================
%  Глава 0 — Введение: кому, зачем, как читать
% ============================================================
\chapter{Зачем эта книжка}
\label{ch:intro}

\section{Кому это}

Этот документ написан для студента, который услышал «модель пространства
состояний» и «MPC» в курсе теории автоматического управления (ТАУ),
открыл код проекта \texttt{SAMURAI} --- и не понял, как одно стыкуется
с другим. Учебник Козлова даёт строгие формулы, но не показывает, как
эти формулы превращаются в команды моторам реального робота. Код проекта
работает и моторы крутятся, но «почему именно эти числа в матрице $\mat{A}$»
не объясняет.

\textbf{Книжка закрывает этот зазор.} Она ведёт читателя по реальной траектории
сценария \texttt{«проехать D метров вперёд»}: от нажатия кнопки в UI до
ШИМ-сигнала в моторе. На каждом шаге появляется ровно та математика,
которая нужна именно здесь, и тут же показывается, в каком файле проекта
она живёт.

\section{Что нужно знать заранее}

\begin{itemize}
  \item \textbf{Матан.} Производные, частные производные, ряд Тейлора первого
        порядка, неопределённый интеграл. Дифференциальные уравнения первого
        порядка с постоянными коэффициентами (типа $\dot{x} = -ax + bu$) ---
        не пугают, но если пугают, то мы их разберём заново.
  \item \textbf{Линейная алгебра.} Что такое матрица, умножение матриц,
        транспонирование, собственные значения. Знать слово «положительно
        определённая» полезно, но не обязательно --- введём.
  \item \textbf{Python на уровне чтения.} Уметь читать сигнатуру функции
        и срез массива \texttt{numpy}. Писать на нём не придётся ---
        главы про реализацию показывают код проекта и комментируют его.
\end{itemize}

Если матан и ЛА давно не открывал --- ничего страшного, нужные вещи
вспомним по ходу.

\section{Чего книжка \emph{не} объясняет}

Чтобы не разрастаться в энциклопедию, документ \textbf{намеренно не покрывает}
следующее:

\begin{itemize}
  \item Электронику и моторы: ШИМ, противо-ЭДС, H-мост. Считаем, что
        мотор --- это «чёрный ящик», который получает команду и крутится
        с некоторой инерцией. Подробнее --- в \filepath{latex\_doc/chapters/14\_pwm\_motor.tex}.
  \item Фильтрацию датчиков: EKF на IMU, слияние акселерометра и одометрии.
        Положение робота уже считаем «измеренным», источник --- глава
        \ref{ch:state}. Подробнее --- в \filepath{latex\_doc/chapters/01\_velocity\_ekf.tex}
        и далее.
  \item SLAM, навигацию A*, Pure Pursuit, патрулирование. Сценарий «D метров
        вперёд» не требует карты и глобального плана --- всё происходит
        в локальной системе координат робота.
\end{itemize}

\section{Как читать}

\textbf{Главы пронумерованы по сценарию использования, а не по сложности
математики.} В каждой главе сначала возникает \emph{вопрос} (например, «откуда
робот знает, на сколько он повернулся?»), потом \emph{интуиция} (на пальцах),
потом \emph{формула}, и в конце --- \emph{ссылка на файл и строку кода},
где это работает.

Если читаешь подряд, к главе~\ref{ch:mpc} ты уже будешь знать, что такое
$\mat{A}$, $\mat{B}$, $\mat{Q}$, $\mat{R}$, и MPC будет уже не «магия из
учебника», а «вот зачем мы их вводили».

Если читаешь нелинейно (только нужное), вот ориентир по типам коробок:

\begin{itemize}
  \item \textbf{Идея} --- интуиция, можно читать без формул.
  \item \textbf{Формула} --- ключевое уравнение, на которое потом ссылаемся.
  \item \textbf{Пример} --- численный пример, можно «прожать» в голове.
  \item \textbf{В коде} --- ссылка на файл/функцию в проекте.
  \item \textbf{Примечание} --- объяснение «в сторону», не обязательное.
  \item \textbf{Что если?} --- учебный «а что будет, если...». Полезно для
        интуиции.
  \item \textbf{Внимание} --- типичная ошибка, реальный подводный камень
        проекта.
\end{itemize}

\section{Что значит «МПС»}

\begin{ideabox}
В этом проекте \textbf{МПС = Модель Пространства Состояний}, по-английски
\emph{state-space model}. Это способ описать поведение робота уравнением
вида
\[
  \vect{x}_{k+1} = \mat{A}_d\,\vect{x}_k + \mat{B}_d\,\vect{u}_k,
\]
где $\vect{x}$ --- вектор состояния (что робот «есть в моменте»: где он,
куда смотрит, как быстро едет), а $\vect{u}$ --- вектор управления (команды
скорости, которые мы ему даём).

Аббревиатура «МПС» \textbf{не} означает «методы принятия решений» (про
Байеса, Гермейера, минимакс) --- если ты искал это, тебе нужна другая
ветка проекта.
\end{ideabox}

\section{Куда смотреть, если книжки мало}

\begin{itemize}
  \item \filepath{latex\_doc/control\_theory/} --- более формальный, сухой
        вариант того же материала. Хорошо для зачёта.
  \item \filepath{latex\_doc/main.tex} --- общий справочник по всей
        математике робота (фильтры, SLAM, кватернионы, A*) --- 15 глав.
  \item \filepath{docs/mps/} --- markdown-доки: \texttt{architecture.md},
        \texttt{api.md}, \texttt{scenario\_forward.md}. Описывают, как
        модуль устроен «снаружи».
  \item \filepath{docs/superpowers/specs/2026-05-05-mps-state-space-design.md}
        --- источник правды по проектным решениям feat/mps.
  \item Учебник: Козлов В.\,Н., Куприянов В.\,Е., Шашихин В.\,Н.
        «Теория автоматического управления», СПбГПУ, 2008.
\end{itemize}

\section{Обозначения}

\begin{center}
\begin{tabular}{ll}
\toprule
$\vect{x}(t) \in \Real^n$    & вектор состояния (непрерывное время)\\
$\vect{x}_k \in \Real^n$     & вектор состояния (дискретное время, шаг $k$)\\
$\vect{u}(t), \vect{u}_k$    & вектор управления, $r$ компонент\\
$\vect{y}(t), \vect{y}_k$    & вектор выхода (что «видит» наблюдатель), $m$ компонент\\
$\mat{A}, \mat{B}$           & матрицы непрерывной модели $\dot{\vect{x}} = \mat{A}\vect{x} + \mat{B}\vect{u}$\\
$\mat{A}_d, \mat{B}_d$       & матрицы дискретной модели $\vect{x}_{k+1} = \mat{A}_d\vect{x}_k + \mat{B}_d\vect{u}_k$\\
$\mat{C}, \mat{D}$           & матрицы выхода: $\vect{y} = \mat{C}\vect{x} + \mat{D}\vect{u}$\\
$\mat{Q}, \mat{R}$           & весовые матрицы критерия качества\\
$\mat{P}_f$                  & терминальный штраф (решение DARE)\\
$\mat{K}$                    & матрица обратной связи регулятора\\
$T_s$                        & период дискретизации (в проекте $0.02$~с)\\
$N$                          & горизонт прогноза MPC\\
$\vect{x}\T$                 & транспонирование\\
$\|\vect{x}\|_{\mat{Q}}^2$   & сокращение для $\vect{x}\T \mat{Q}\, \vect{x}$ (квадратичная форма)\\
$\bar{\vect{x}}, \bar{\vect{u}}$ & опорная точка линеаризации\\
$\Delta\vect{x}$             & отклонение от опорной точки: $\vect{x} - \bar{\vect{x}}$\\
\bottomrule
\end{tabular}
\end{center}

Переменные с тильдой (\texttt{\~{}x}) встречаются в коде --- это та же
$\vect{x}$, просто иначе оформлено в идентификаторах.

\medskip
Готов? Тогда поехали ---  смотрим, что вообще делает робот.

% ============================================================
%  Глава 1 — Что робот вообще делает
% ============================================================
\chapter{Робот глазами пользователя}
\label{ch:user_view}

\section{Сценарий «D метров вперёд»}

Открываем веб-дашборд (\filepath{compute\_node/frontend/}), идём на страницу
\texttt{/mps}. Видим окружность с маркером, поля «\textit{дистанция $D$}»
и «\textit{целевая скорость $v_\text{target}$}», кнопку «Run on Robot».

\begin{enumerate}
  \item Двигаем маркер на окружности --- выбираем \textbf{целевой курс}
        $\varphi$ (на сколько градусов отвернуть от «прямо вперёд»).
  \item Вводим $D$ и $v_\text{target}$.
  \item Жмём «Run on Robot».
\end{enumerate}

Робот делает три вещи:

\begin{enumerate}
  \item \textbf{Поворачивается на месте} на угол $\varphi$. Едет на нуле,
        крутится только разворот.
  \item \textbf{Проезжает $D$ метров} с целевой скоростью $v_\text{target}$,
        удерживая курс $\varphi$.
  \item \textbf{Останавливается} и публикует результат: сколько проехал,
        какой был перерегулирующий выброс, как долго «успокаивался».
\end{enumerate}

Всё. Один сценарий, конечный. Никакой карты, никакого глобального плана.

\begin{ideabox}
Книжка построена вокруг \emph{этого} сценария. Каждый следующий вопрос
--- «а как робот\dots» --- разбираем именно для этого сценария. Простота
сценария --- наше учебное преимущество: можно вглядеться в каждую формулу
без отвлечения на другие задачи.
\end{ideabox}

\section{Что нужно роботу, чтобы это сделать}

Распакуем то, что мы только что увидели. У робота должна быть способность
ответить на пять вопросов.

\begin{enumerate}
  \item[\textbf{(Q1)}] \textbf{Где я?}
        Чтобы понять, что проехал $D$ метров, надо знать, на сколько
        я уже сдвинулся. Это \emph{состояние позиции}.
        Об этом --- глава~\ref{ch:state}.

  \item[\textbf{(Q2)}] \textbf{Куда я смотрю?}
        Чтобы повернуться к $\varphi$ и не уехать боком, надо знать
        текущий курс $\theta$. Это \emph{состояние курса}.
        Тоже глава~\ref{ch:state}.

  \item[\textbf{(Q3)}] \textbf{Как мне сдвинуться/повернуться, если я скажу
        моторам «езжай»?}
        Если послать команду $u_v = 0.2$~м/с, робот не сразу станет ехать
        $0.2$~м/с --- моторы инерционны. Это \emph{динамика}.
        Об этом --- главы~\ref{ch:kinematics}, \ref{ch:dynamics}.

  \item[\textbf{(Q4)}] \textbf{Какую команду подать, чтобы достичь цели?}
        Не любой набор команд приведёт к цели хорошо. Нужен \emph{регулятор}.
        Об этом --- главы~\ref{ch:controllers}, \ref{ch:lqr}, \ref{ch:mpc}.

  \item[\textbf{(Q5)}] \textbf{Когда остановиться и как сказать «готово»?}
        Это \emph{логика сценария}, конечный автомат и подсчёт метрик.
        Об этом --- главы~\ref{ch:scenario}, \ref{ch:turn_phase}.
\end{enumerate}

\section{Три слоя проекта}

Робот --- не один компьютер, а три. Они общаются через сеть.

\begin{center}
\begin{tikzpicture}[
  node distance=8mm,
  layer/.style={
    draw, rounded corners=3pt, thick,
    align=center, font=\small,
    minimum width=4.6cm, minimum height=1.7cm,
  },
  flow/.style={-{Stealth[length=2.5mm]}, thick}
]
  \node[layer, fill=blue!8]   (ui)  {\textbf{Browser / React}\\
                                     UI \texttt{/mps}, MatrixEditor,\\
                                     Target picker};
  \node[layer, fill=green!8, right=of ui]  (cn) {\textbf{Compute Node (ноут)}\\
                                     FastAPI, MQTT-broker,\\
                                     idealised simulator};
  \node[layer, fill=orange!8, right=of cn] (pi) {\textbf{Raspberry Pi 4}\\
                                     \texttt{mps\_node.py}, FSM,\\
                                     MPC, motor driver};
  \draw[flow] (ui) -- node[above, font=\scriptsize] {REST + WS} (cn);
  \draw[flow] (cn) -- node[above, font=\scriptsize] {MQTT} (pi);
\end{tikzpicture}
\end{center}

Что где живёт:

\begin{itemize}
  \item \textbf{Browser} (React) --- то, что видит человек: пикер курса,
        поля распределения, графики прогона. Жесть математики тут
        минимум, но есть --- глава~\ref{ch:picker} разберёт.
  \item \textbf{Compute Node} (FastAPI на ноутбуке) --- посредник: принимает
        команды от UI, шлёт их роботу по MQTT, сохраняет историю прогонов.
        Заодно умеет \emph{симулировать} прогон без физического робота
        --- это идеальный $\vect{x}_{k+1} = \mat{A}_d\vect{x}_k + \mat{B}_d\vect{u}_k$
        без шумов. Полезно для UI и юнит-тестов.
  \item \textbf{Pi} (Python, MQTT) --- собственно мозг робота. Подписан
        на топик \texttt{mps/scenario/run}; запускает тик-цикл $50$~Гц,
        внутри которого \emph{вся боевая математика} --- считывание
        одометрии, шаг MPC, выдача команды моторам.
\end{itemize}

\begin{codebox}[title={Где что лежит}]
\begin{tabular}{ll}
\filepath{compute\_node/frontend/src/pages/MpsPage.tsx} & UI страница\\
\filepath{compute\_node/frontend/src/lib/targetAngle.ts} & математика пикера\\
\filepath{compute\_node/dashboard/routers/mps.py}        & REST + WS\\
\filepath{compute\_node/mps\_runner.py}                  & идеальный симулятор\\
\filepath{pi\_nodes/nodes/mps\_node.py}                  & оркестратор на Pi\\
\filepath{pi\_nodes/control/state\_space\_model.py}      & plant ($\mat{A}_d, \mat{B}_d$)\\
\filepath{pi\_nodes/control/mpc\_controller.py}          & MPC ($\mat{Q}, \mat{R}$, $\mat{K}$, QP)\\
\filepath{config.yaml}                                   & секция \texttt{mps:} с матрицами\\
\filepath{matlab/main.m}                                 & оффлайн-синтез матриц\\
\end{tabular}
\end{codebox}

\section{Граничные правила}

Перед тем как нырять в математику, зафиксируем правила игры, которые
не выводятся, а просто заданы дизайном модуля. Если ты увидишь в коде
проверку, которая кажется лишней --- скорее всего, она охраняет одно
из этих правил.

\begin{notebox}
\begin{enumerate}
  \item MPC активен \textbf{только в FSM-состоянии} \texttt{DRIVE\_FORWARD\_MPS}.
        В любом другом (IDLE, SEARCHING, GRABBING, ...) команды моторам
        идут от \emph{другого} контроллера (старый PID), а MPC ничего не
        делает.
  \item Сценарий \textbf{финитный}: достижение цели, тайм-аут, abort или
        ошибка --- любое из четырёх событий завершает прогон, и FSM
        возвращается в IDLE.
  \item Между прогонами можно \textbf{пересобрать} матрицы (MQTT-команда
        \texttt{mps/matrices/set}) --- это и есть «учебный пульт».
        Во время прогона матрицы не меняются.
  \item В состоянии \texttt{DRIVE\_FORWARD\_MPS} \textbf{cmd\_vel} (команда
        моторам) приходит \emph{только} от \texttt{mps\_node}. Голосовые
        команды, детекции мяча и джойстик игнорируются.
  \item Симулятор \textbf{идеальный}: $\vect{x}_{k+1} = \mat{A}_d\vect{x}_k +
        \mat{B}_d\vect{u}_k$ без шумов, без \emph{motor lag}, без проскальзывания.
        Робот --- неидеальный, разница между симом и реальным прогоном
        --- это и есть «ошибка моделирования».
\end{enumerate}
\end{notebox}

\section{Что дальше}

Дальше мы аккуратно ответим на пять вопросов (Q1)--(Q5) по очереди. Сначала
$\to$ кинематика (глава~\ref{ch:kinematics}): что значит «робот едет вперёд»
формально. Потом --- откуда состояние берётся (глава~\ref{ch:state}). Затем
динамика, линеаризация, регуляторы, и наконец --- сам сценарий.

Идём.

% ============================================================
%  Глава 2 — Дифференциальная кинематика
% ============================================================
\chapter{Дифференциальная кинематика}
\label{ch:kinematics}

\section{Вопрос}

Робот --- гусеничное шасси с \textbf{двумя задними моторами} (передние
физически отсутствуют, см. \filepath{hardware\_presets/samurai\_default.yaml}).
Левый мотор крутит левую гусеницу, правый --- правую.

\begin{ideabox}
Мы хотим сказать роботу: «езжай со скоростью $0.2$~м/с и поворачивай
со скоростью $0.5$~рад/с». Внутри у него --- две гусеницы, которые
крутятся каждая со своей скоростью. Откуда взять скорости гусениц
из «езжай вперёд + поворачивай»?
\end{ideabox}

\section{Геометрия}

Введём систему координат, связанную с роботом. Робот ---  это «жёсткий
прямоугольник» шириной $B$ метров (расстояние между серединами левой и
правой гусениц, \texttt{wheel\_base}). Ось $x_b$ робота смотрит вперёд,
ось $y_b$ --- влево.

\begin{center}
\begin{tikzpicture}[scale=1.1, every node/.style={font=\small}]
  % Корпус робота
  \draw[thick, fill=gray!10, rounded corners=2pt]
    (-1.0, -0.7) rectangle (1.0, 0.7);
  % Левая и правая гусеницы (полоски)
  \fill[gray!50] (-1.05, 0.45) rectangle (1.05, 0.7);
  \fill[gray!50] (-1.05,-0.7)  rectangle (1.05,-0.45);
  % Стрелка вперёд
  \draw[-{Stealth[length=3mm]}, thick, blue!60!black] (0,0) -- (1.6,0)
       node[right]{\(x_b\) (вперёд)};
  % Ось y_b
  \draw[-{Stealth[length=3mm]}, thick, blue!60!black] (0,0) -- (0,1.4)
       node[above]{\(y_b\) (влево)};
  % Скорости гусениц
  \draw[-{Stealth[length=2.5mm]}, very thick, red!70!black]
       (-0.8, 0.575) -- (0.5, 0.575)
       node[right]{\(v_L\)};
  \draw[-{Stealth[length=2.5mm]}, very thick, red!70!black]
       (-0.8,-0.575) -- (0.2,-0.575)
       node[right]{\(v_R\)};
  % Ширина базы
  \draw[<->, thick] (-1.25, 0.575) -- (-1.25, -0.575)
       node[midway, left]{\(B\)};
\end{tikzpicture}
\end{center}

Пусть гусеницы крутятся со скоростями $v_L$ (левая) и $v_R$ (правая)
--- это \textbf{линейные} скорости центров гусениц в продольном
направлении, м/с.

\section{Команды робота: $u_v$ и $u_\omega$}

Команды роботу --- не «крути левую и правую», а более удобная пара:
\[
  \vect{u} = \begin{pmatrix} u_v \\ u_\omega \end{pmatrix},
  \qquad
  u_v \text{ -- линейная скорость [м/с]},
  \qquad
  u_\omega \text{ -- угловая скорость [рад/с]}.
\]

\begin{itemize}
  \item $u_v > 0$ --- ехать вперёд, $u_v < 0$ --- назад.
  \item $u_\omega > 0$ --- поворот \emph{налево} (CCW против часовой,
        стандарт ROS и наш проект).
  \item $u_v = 0$, $u_\omega \ne 0$ --- разворот на месте.
  \item $u_v \ne 0$, $u_\omega \ne 0$ --- движение по дуге.
\end{itemize}

\begin{formulabox}[title={Прямая дифференциальная кинематика}]
Из геометрии: центр шасси едет со скоростью $u_v$, а каждая гусеница ---
на расстоянии $\tfrac{B}{2}$ от центра, поэтому добавочно скользит на
$\pm\tfrac{B}{2}\,u_\omega$ (знак зависит от того, левая она или правая).

\begin{equation}
  \boxed{\;v_L = u_v + \tfrac{B}{2}\,u_\omega,\qquad
         v_R = u_v - \tfrac{B}{2}\,u_\omega.\;}
  \label{eq:diff_forward}
\end{equation}
\end{formulabox}

\begin{notebox}
\textbf{Откуда знаки?} При повороте налево ($u_\omega > 0$) робот «закидывает
зад» вправо, его левая гусеница должна крутиться \emph{быстрее} центра,
а правая --- медленнее. Поэтому к $u_v$ добавляем $+\tfrac{B}{2}u_\omega$
для левой и вычитаем для правой.
\end{notebox}

\section{Обратная задача: гусеницы $\to$ команды}

Часто нужна обратная задача --- по измеренным скоростям гусениц найти
$u_v$ и $u_\omega$. Складываем и вычитаем \eqref{eq:diff_forward}:

\begin{equation}
  u_v = \frac{v_L + v_R}{2},
  \qquad
  u_\omega = \frac{v_L - v_R}{B}.
  \label{eq:diff_inverse}
\end{equation}

\begin{examplebox}[title={Численный пример}]
$B = 0.17$~м. Хотим $u_v = 0.2$~м/с, $u_\omega = 0.5$~рад/с.

\begin{align*}
  v_L &= 0.2 + \tfrac{0.17}{2}\cdot 0.5 = 0.2 + 0.0425 = 0.2425 \text{ м/с},\\
  v_R &= 0.2 - 0.0425 = 0.1575 \text{ м/с}.
\end{align*}

Левая гусеница на 30\% быстрее правой $\to$ робот поворачивает налево.
\end{examplebox}

\section{Положение и курс в мировой системе}

Команды $u_v$, $u_\omega$ --- в \emph{связанной} с роботом системе. Но
позиция $(p_x, p_y)$ и курс $\theta$ --- в \emph{мировой}. Чтобы
получить их, нужно вспомнить, куда робот смотрит.

\begin{ideabox}
Если робот смотрит под углом $\theta$ к мировой оси $x$ и движется вперёд
со скоростью $v$, то проекции этой скорости на мировые оси --- $v\cos\theta$
и $v\sin\theta$. Курс $\theta$ меняется со скоростью $\omega$.
\end{ideabox}

\begin{formulabox}[title={Кинематика робота в мировой СК}]
\begin{equation}
  \begin{cases}
    \dot{p}_x = v\,\cos\theta,\\[2pt]
    \dot{p}_y = v\,\sin\theta,\\[2pt]
    \dot{\theta} = \omega.
  \end{cases}
  \label{eq:world_kinematics}
\end{equation}
\end{formulabox}

Это \textbf{первое нелинейное звено} модели. Нелинейность здесь --- из-за
$\cos\theta$ и $\sin\theta$, зависящих от состояния.

\section{Где это в коде}

\begin{codebox}[title={\filepath{pi\_nodes/nodes/motor\_node.py}}]
В микшере моторов \eqref{eq:diff_forward} применяется к \texttt{cmd\_vel},
полученному из MPC. Дальше скорости гусениц переводятся в ШИМ через
калибровочные коэффициенты (PCA9685, каналы M1 ch10--11, M2 ch8--9).
Подробнее --- глава \filepath{latex\_doc/chapters/14\_pwm\_motor.tex}.
\end{codebox}

\begin{codebox}[title={\filepath{config.yaml}, секция \texttt{simulator:}}]
\begin{verbatim}
simulator:
  wheel_base:  0.17    # B [м]
  max_linear:  0.30    # |u_v|_max [м/с]
  max_angular: 2.00    # |u_omega|_max [рад/с]
\end{verbatim}
\end{codebox}

\section{Что мы получили}

\begin{itemize}
  \item Команда роботу --- это пара $(u_v, u_\omega)$, размерность $r = 2$.
  \item Мгновенно «исполнить» команду робот не может --- сначала надо
        перевести её в скорости гусениц (тривиально, \eqref{eq:diff_forward}),
        а потом --- в ШИМ (см. отдельную главу). \emph{Это всё ещё чистая
        кинематика без инерции}.
  \item Перевод в мировые координаты --- через \eqref{eq:world_kinematics}.
        Здесь появилась первая нелинейность модели.
\end{itemize}

В следующей главе разберёмся, \emph{откуда} робот узнаёт свои $p_x, p_y, \theta$
(и можно ли им верить).

% ============================================================
%  Глава 3 — Откуда берётся состояние
% ============================================================
\chapter{Откуда берётся состояние}
\label{ch:state}

\section{Вопрос}

В \eqref{eq:world_kinematics} мы писали $\dot{p}_x = v\cos\theta$. Это
дифференциальное уравнение --- то, как состояние \emph{меняется во времени}.
Но MPC, который мы будем запускать $50$~раз в секунду, ему на каждом тике
нужно $\vect{x}_k$ --- \emph{измеренное} состояние \emph{сейчас}.

Откуда оно? Робот же не «знает» своих координат --- их нужно посчитать
из сигналов реальных датчиков.

\begin{ideabox}
Это разделение --- очень важное:
\begin{itemize}
  \item \emph{Модель} говорит, как состояние должно меняться (формула).
  \item \emph{Оценка состояния} (state estimation) даёт текущее значение
        состояния по показаниям датчиков (число).
\end{itemize}
В этой главе нас интересует \emph{оценка}. Модель будет дальше.
\end{ideabox}

\section{Состояние в МПС-сценарии}

В сценарии «D метров вперёд» состояние --- это пять чисел:

\begin{formulabox}[title={Канонический МПС-вектор состояния}]
\begin{equation}
  \vect{x} \;=\; \begin{pmatrix} s \\ v \\ \theta \\ \omega \\ e_\text{int} \end{pmatrix}
  \;\in\; \Real^5
  \label{eq:mps_state}
\end{equation}

\begin{tabular}{lll}
\toprule
$s$           & пройденная дистанция         & м\\
$v$           & продольная скорость           & м/с\\
$\theta$      & курс (yaw)                    & рад\\
$\omega$      & угловая скорость              & рад/с\\
$e_\text{int}$& интеграл ошибки скорости      & м ($\approx v\cdot t$)\\
\bottomrule
\end{tabular}
\end{formulabox}

\begin{notebox}
\textbf{Почему именно так, а не $(p_x, p_y, \theta, v, \omega)$?}

В общей литературе по дифференциальным роботам состояние ---
$(p_x, p_y, \theta, v, \omega)$ (две декартовы координаты + курс +
две скорости). И в проекте есть «легаси»-форма с такими пятью переменными
(\filepath{pi\_nodes/control/state\_space\_model.py}: переменные
\texttt{px, py, theta, v, omega}). Но для сценария «D метров вперёд»
эта форма \emph{избыточна}: при правильной работе $p_y \equiv 0$ (робот
едет прямо), а $p_x$ совпадает с $s$.

Поэтому в МПС-сценарии мы переходим к более компактной форме
\eqref{eq:mps_state}: одна продольная координата $s$ вместо двух,
плюс \emph{интегральный член} $e_\text{int}$, который ловит установившуюся
ошибку по скорости (без него робот хронически не добирает $v_\text{target}$).
\end{notebox}

\begin{warnbox}
В коде live ОБЕ модели существуют:
\begin{itemize}
  \item Легаси $(p_x, p_y, \theta, v, \omega)$ --- секция \texttt{control.}
        в \filepath{config.yaml}, используется старым PID-контуром.
  \item Каноническая $(s, v, \theta, \omega, e_\text{int})$ --- секция
        \texttt{mps:}, используется в MPS-сценарии.
\end{itemize}
Это \emph{не баг} --- это разделение «учебной» подсистемы и боевой.
Просто когда читаешь код, всегда смотри, \emph{какая} секция
(\texttt{control.*} или \texttt{mps.*}) грузит матрицы.

Историческая ловушка: \texttt{mps\_node} в коммите \texttt{8521742}
случайно строился из \texttt{control.*} и трактовал свой
$x[1] = v$ как $p_y$, после чего MPC доворачивал робота, пытаясь
свести $p_y$ к $v_\text{target}$. Робот ездил по кругу.
\end{warnbox}

\section{Источники измерений}

Каждая компонента \eqref{eq:mps_state} приходит не «магически», а из
конкретного сенсора через конкретный фильтр.

\begin{center}
\begin{tabular}{llp{6cm}}
\toprule
\textbf{Компонента} & \textbf{Сенсор} & \textbf{Преобразование}\\
\midrule
$s, v$    & энкодеры моторов   & интегрирование числа тиков $\to$ путь;
                                 \texttt{velocity\_ekf}\\
$\theta, \omega$ & MPU-6050 IMU (гироскоп) & \texttt{ekf\_imu}: фильтр Калмана
                                            1D на $\theta$\\
$e_\text{int}$ & --- & накапливается \emph{самим} \texttt{mps\_node} как
                       $\sum (v - v_\text{target})\,\Delta t$\\
\bottomrule
\end{tabular}
\end{center}

Подробное описание фильтров --- в общем справочнике
\filepath{latex\_doc/chapters/01\_velocity\_ekf.tex},
\filepath{latex\_doc/chapters/02\_imu\_ekf.tex}.

\subsection{Одометрия по гусеницам}

Энкодеры считают, сколько раз мотор провернулся за последний тик.
Зная диаметр приводного колеса и передаточное число, это превращается
в пройденный гусеницей путь $\Delta s_L, \Delta s_R$. Скорости и
позиция центра:

\begin{equation}
  v = \frac{\Delta s_L + \Delta s_R}{2\,\Delta t},
  \qquad
  \omega_\text{odom} = \frac{\Delta s_L - \Delta s_R}{B\,\Delta t},
  \qquad
  s_{k+1} = s_k + v\,\Delta t.
\end{equation}

Эти формулы --- та же \eqref{eq:diff_inverse}, что в главе~\ref{ch:kinematics},
только применённая к \emph{измеренным} $\Delta s$, а не к команде.

\subsection{Курс по IMU}

Гироскоп MPU-6050 даёт измеренную угловую скорость $\omega_\text{imu}(t)$.
Интегрируя её, получим курс $\theta(t) = \int \omega_\text{imu}\,dt$.
Проблема: гироскоп дрейфует. За минуту накопится несколько градусов
ошибки даже на стоящем роботе.

Для коротких сценариев ($D \le 5$~м $\Rightarrow$ $t \le 25$~с при
$v_\text{target} = 0.2$~м/с) дрейф пренебрежимо мал, и мы используем
прямое интегрирование с лёгкой EKF-сглаживанием. Для длинных сессий
проект делает дополнительную привязку через \texttt{position\_fusion}
(акселерометр + одометрия), но в МПС-сценарии это не критично.

\section{Слияние: \texttt{position\_fusion}}

Все источники сводятся в одну ноду \texttt{position\_fusion}, которая
публикует MQTT-сообщение \texttt{odom} с полями \texttt{x, vx, theta, vz}.
Это и есть «питающие» измерения для \texttt{mps\_node}.

\begin{codebox}[title={\filepath{pi\_nodes/nodes/mps\_node.py:342}}]
\begin{lstlisting}[language=Python]
def _on_odom(self, topic: str, payload):
    # samurai/{id}/odom — motor_node публикует:
    #   x в САНТИМЕТРАХ, vx в м/с, theta в рад, vz в рад/с.
    s = float(payload.get('x', 0.0)) / 100.0   # см -> м
    v = float(payload.get('vx', 0.0))
    theta = float(payload.get('theta', 0.0))
    omega = float(payload.get('vz', 0.0))
    with self._lock:
        self._x_meas = np.array(
            [s, v, theta, omega, self._x_meas[_EINT]]
        )
        self._x_meas_ts = time.time()
\end{lstlisting}
\end{codebox}

\begin{warnbox}
\textbf{Подводный камень --- единицы.} \texttt{motor\_node} публикует
позицию в \emph{сантиметрах}, а МПС-модель работает в \emph{метрах}.
В коммите \texttt{5f54c54} баг был такой: эвристика «если $|s| > 20$,
то это уже метры, не конвертируй» оставляла первые 20~см не
сконвертированными --- и MPC видел фантомную ошибку позиции $\times 100$,
прижимая $u$ к насыщению $u_\text{max}$. Исправлено безусловной
конвертацией.

Урок: в любой межкомпонентной границе \emph{явно фиксируй единицы}.
\end{warnbox}

\section{Старт сценария: «обнуление» начала}

Когда пользователь жмёт «Run», \texttt{mps\_node} запоминает \emph{текущие}
$s$ и $\theta$ как $s_\text{start}, \theta_\text{start}$ (см. \texttt{\_RunState}).
Дальше состояние для MPC считается \textbf{относительно старта}:

\begin{equation}
  x[\,s\,]_\text{rel}     = s - s_\text{start},
  \qquad
  x[\,\theta\,]_\text{rel} = \text{wrap}_{[-\pi, \pi]}(\theta - \theta_\text{start}).
  \label{eq:relative_state}
\end{equation}

\begin{ideabox}
Зачем? Сценарий говорит: «проехать $D$ метров \emph{отсюда}». Если бы
мы кормили в MPC \emph{абсолютную} $s$, то reference $s_\text{ref}$
тоже надо было бы делать абсолютным --- а робот мог стартовать с
$s=4.7$~м после предыдущего прогона. Гораздо проще обнулить отсчёт
в момент старта.

То же с курсом: $\theta_\text{ref} = \varphi$ означает «отвернуть на
$\varphi$ от \emph{текущего} направления», а не от мирового севера.
\end{ideabox}

Функция \texttt{\_normalize\_angle} оборачивает разность в $[-\pi, \pi]$,
чтобы поворот на $-179^\circ$ не превратился в $+181^\circ$ (что бы
сломало MPC, который квадратично штрафует ошибку курса).

\section{Что мы получили}

\begin{itemize}
  \item Состояние МПС-сценария --- пять чисел $(s, v, \theta, \omega, e_\text{int})$.
  \item Первые четыре приходят из MQTT-топика \texttt{odom} с Pi-ноды
        \texttt{position\_fusion}, конвертируясь см$\to$м.
  \item Пятый ($e_\text{int}$) живёт внутри самого \texttt{mps\_node}
        --- это интеграл ошибки скорости, нужен для устойчивого
        отслеживания $v_\text{target}$.
  \item На старте прогона $s$ и $\theta$ \emph{обнуляются}: дальнейшие
        вычисления идут в относительных координатах.
\end{itemize}

Теперь у нас есть и состояние, и формула \eqref{eq:world_kinematics}, как
оно меняется кинематически. Но кинематики мало: моторы инерционны,
команда отрабатывается не мгновенно. Это --- следующая глава.

% ============================================================
%  Глава 4 — Динамика: инерция моторов
% ============================================================
\chapter{Динамика: моторы не отрабатывают мгновенно}
\label{ch:dynamics}

\section{Вопрос}

В предыдущей главе мы кормили в кинематику $v$ и $\omega$ так, будто
\emph{что задали --- то и едет}. На самом деле это не так. Если в
момент $t=0$ сказать роботу «езжай $0.2$~м/с» (раньше стоял), он
\emph{не сразу} достигнет $0.2$~м/с --- сначала будет разгоняться.

\begin{ideabox}
Между \emph{командой} $u_v$ и \emph{фактической} скоростью $v$ есть
зазор: моторам нужно время, чтобы раскрутить инерцию, преодолеть
противо-ЭДС и трение. Длительность этого зазора --- постоянная времени
$\tau_v$.

Это \textbf{динамика} робота. Без неё модель неполна, и MPC,
обученный на одной только кинематике, будет систематически ошибаться.
\end{ideabox}

\section{Апериодическое звено первого порядка}

Простейшая модель, которая адекватно описывает разгон моторов под
такой нагрузкой --- \textbf{апериодическое звено первого порядка}
(в литературе также «инерционное звено», PT$_1$):

\begin{formulabox}[title={Динамика линейной скорости}]
\begin{equation}
  \dot{v} = \frac{1}{\tau_v}\,(u_v - v).
  \label{eq:motor_lin_dyn}
\end{equation}
\end{formulabox}

Если $u_v$ постоянно, решение этого уравнения --- классическая экспонента:

\begin{equation}
  v(t) = u_v - (u_v - v_0)\,e^{-t/\tau_v},
  \label{eq:exp_response}
\end{equation}

где $v_0 = v(0)$. За время $t = \tau_v$ робот покрывает $\approx 63\%$
пути от $v_0$ до $u_v$; за $t = 3\tau_v$ --- $\approx 95\%$.

\begin{ideabox}
Можно прочитать \eqref{eq:motor_lin_dyn} «по-человечески»: \emph{скорость
изменения скорости пропорциональна тому, как сильно она отстаёт от
команды}. Чем больше отставание --- тем сильнее «толкаем»; когда $v$
догнала $u_v$ --- толкать перестаём.

Аналогия: вода, текущая в бак, где скорость наполнения пропорциональна
разности уровней. RC-цепь в электронике --- та же самая динамика, потому
её и называют «RC-звено».
\end{ideabox}

\begin{examplebox}[title={Численный пример}]
Пусть $\tau_v = 0.15$~с (типичное значение для шасси SAMURAI).
Робот стоит ($v_0 = 0$), команда $u_v = 0.2$~м/с.

\begin{tabular}{rrr}
\toprule
$t$, с & $v$, м/с & доля от $u_v$\\
\midrule
$0.05$ & $0.2 \cdot (1 - e^{-1/3}) \approx 0.057$ & 28.3\%\\
$0.15$ & $0.2 \cdot (1 - e^{-1}) \approx 0.126$ & 63.2\%\\
$0.30$ & $0.2 \cdot (1 - e^{-2}) \approx 0.173$ & 86.5\%\\
$0.45$ & $0.2 \cdot (1 - e^{-3}) \approx 0.190$ & 95.0\%\\
\bottomrule
\end{tabular}

За $T_s = 0.05$~с (один тик MPC) робот успевает разогнаться только
до 28\% желаемого. Это и есть то, что MPC обязан знать.
\end{examplebox}

\section{Аналогично --- угловая скорость}

\begin{formulabox}[title={Динамика угловой скорости}]
\begin{equation}
  \dot{\omega} = \frac{1}{\tau_\omega}\,(u_\omega - \omega).
  \label{eq:motor_ang_dyn}
\end{equation}
\end{formulabox}

Постоянная $\tau_\omega \approx 0.10$~с заметно \emph{меньше} $\tau_v$:
крутить шасси на месте проще, чем разгонять центр массы.

\section{Полная нелинейная ОДУ-модель}

Соберём кинематику \eqref{eq:world_kinematics} и динамику
\eqref{eq:motor_lin_dyn}, \eqref{eq:motor_ang_dyn} в один вектор.
Возьмём «физическое» (легаси) состояние из главы~\ref{ch:state}:

\[
  \vect{x}(t) = \big(p_x,\; p_y,\; \theta,\; v,\; \omega\big)\T \in \Real^5,
  \qquad
  \vect{u}(t) = (u_v,\; u_\omega)\T \in \Real^2.
\]

\begin{formulabox}[title={Нелинейная ОДУ-модель робота}]
\begin{equation}
  \boxed{\;
    \dot{\vect{x}}(t) = \vect{f}(\vect{x}, \vect{u})
    \;=\;
    \begin{pmatrix}
      v\cos\theta \\[3pt]
      v\sin\theta \\[3pt]
      \omega \\[3pt]
      \frac{1}{\tau_v}\,(u_v - v) \\[5pt]
      \frac{1}{\tau_\omega}\,(u_\omega - \omega)
    \end{pmatrix}.\;}
  \label{eq:nonlinear_ode}
\end{equation}
\end{formulabox}

Здесь:
\begin{itemize}
  \item Первые три строки --- \emph{кинематика}. Нелинейность через $\cos, \sin$.
  \item Последние две --- \emph{динамика моторов}. \emph{Линейная}: правая
        часть линейна по $v, \omega, u_v, u_\omega$.
\end{itemize}

\textbf{Вся нелинейность сосредоточена в куске «как $v$ и $\theta$
порождают движение в мире»}. Это важно: при линеаризации (следующая
глава) именно этот кусок упростится.

\section{Откуда берутся $\tau_v$ и $\tau_\omega$}

Постоянные времени --- \emph{физические} параметры, оценённые
экспериментально. В \filepath{tests/test\_motor\_lag.py} есть тест, который
гонит робот на серии импульсов $u_v$ и подбирает $\tau_v$
по методу наименьших квадратов из закона \eqref{eq:exp_response}.
В \filepath{config.yaml} они зашиты как:

\begin{codebox}[title={\filepath{config.yaml}, секция \texttt{control.plant}}]
\begin{verbatim}
control:
  plant:
    Ts:     0.05     # период дискретизации, с
    v0:     0.20     # опорная скорость для линеаризации, м/с
    tau_v:  0.15     # постоянная времени линейной скорости, с
    tau_w:  0.10     # постоянная времени угловой скорости, с
\end{verbatim}
\end{codebox}

\begin{whatifbox}
\textbf{Что если у моторов больше инерция?} Тогда $\tau_v$ увеличится,
и в дискретной модели $1/\tau_v$ уменьшится --- собственное значение
$\lambda = -1/\tau_v$ ближе к нулю. В дискретном виде $|\tilde{\lambda}|
= e^{-T_s/\tau_v}$ ближе к единице --- мотор «медленнее забывает» прошлую
команду. MPC должен это учесть, выбирая более «нежную» $\mat{R}$.

Если поставить в config \emph{неправильно большое} $\tau_v$, ничего
не сломается --- модель просто будет «думать, что моторы вялые». MPC
будет давать более резкие команды (чтоб компенсировать выдуманную
инерцию) --- робот будет дёргаться. Хороший пример «учебной катастрофы»
из задачи курсовой.
\end{whatifbox}

\section{Где это в коде}

\begin{codebox}[title={\filepath{pi\_nodes/control/state\_space\_model.py:29}}]
\begin{lstlisting}[language=Python]
def _continuous_AB(v0: float, tau_v: float, tau_w: float):
    """Build linearised A, B around forward motion at v0."""
    A = np.array([
        [0.0, 0.0, 0.0,        1.0,         0.0       ],
        [0.0, 0.0, v0,         0.0,         0.0       ],
        [0.0, 0.0, 0.0,        0.0,         1.0       ],
        [0.0, 0.0, 0.0,       -1.0/tau_v,   0.0       ],
        [0.0, 0.0, 0.0,        0.0,        -1.0/tau_w ],
    ])
    B = np.array([
        [0.0,        0.0       ],
        [0.0,        0.0       ],
        [0.0,        0.0       ],
        [1.0/tau_v,  0.0       ],
        [0.0,        1.0/tau_w ],
    ])
    return A, B
\end{lstlisting}
\end{codebox}

Это уже \emph{линеаризованная} версия модели --- но мы её ещё не вывели.
Чтобы её получить, нужно сделать линеаризацию \eqref{eq:nonlinear_ode} вокруг
опорной точки. Этим займёмся в следующей главе.

\section{Что мы получили}

\begin{itemize}
  \item Моторы --- \emph{инерционные}: команда отрабатывается за $\tau_v
        \sim 0.15$~с, $\tau_\omega \sim 0.10$~с.
  \item Динамика моторов --- линейная: $\dot{v} = (u_v - v)/\tau_v$ и
        аналогично для $\omega$.
  \item Полная модель робота \eqref{eq:nonlinear_ode} --- пятимерная,
        нелинейная (только в кинематической части), с двумя управлениями.
  \item Постоянные времени --- из тестов на реальном железе, сидят
        в \filepath{config.yaml}.
\end{itemize}

Дальше эту нелинейную систему надо превратить в линейную (для MPC и
ЛКР это базовая необходимость). Идёмм.

% ============================================================
%  Глава 5 — Линеаризация и дискретизация
% ============================================================
\chapter{Из нелинейной ОДУ в матрицы $\mat{A}_d, \mat{B}_d$}
\label{ch:linearization}

\section{Вопрос}

Нелинейная модель \eqref{eq:nonlinear_ode} --- честная, но \textbf{нерешаемая
аналитически} и плохо подходящая для оптимизаторов. ЛКР и MPC, которые
мы будем строить дальше, работают \emph{только с линейными моделями}
в дискретном времени:

\[
  \vect{x}_{k+1} = \mat{A}_d\,\vect{x}_k + \mat{B}_d\,\vect{u}_k.
\]

Надо превратить \eqref{eq:nonlinear_ode} в эту форму.

\begin{ideabox}
Две операции, обе из учебника Козлова:
\begin{enumerate}
  \item \textbf{Линеаризация.} Разложение в ряд Тейлора первого порядка
        вокруг \emph{опорной точки} (равномерное движение вперёд).
        Получаем \emph{непрерывные} матрицы $\mat{A}, \mat{B}$.
  \item \textbf{Дискретизация.} Переход от $\dot{\vect{x}} = \mat{A}\vect{x}
        + \mat{B}\vect{u}$ к $\vect{x}_{k+1} = \mat{A}_d\vect{x}_k +
        \mat{B}_d\vect{u}_k$ через ZOH (zero-order hold).
\end{enumerate}
В этой главе пройдём по обеим.
\end{ideabox}

\section{Выбор опорной точки}

\textbf{Опорная точка} (operating point, англ.\,\emph{trim point}) --- это
тот режим, вокруг которого мы линеаризуем. Выберем его так, чтобы он
был \emph{реальной целью} сценария: робот едет вперёд с постоянной
скоростью $v_0$, ничего не вращая.

\begin{equation}
  \bar{\vect{x}} = (\bar{p}_x,\; \bar{p}_y,\; \underbrace{0}_{\bar\theta},\;
                    v_0,\; \underbrace{0}_{\bar\omega})\T,
  \qquad
  \bar{\vect{u}} = (v_0,\; 0)\T.
  \label{eq:op_point}
\end{equation}

Подставим $\bar{\vect{x}}, \bar{\vect{u}}$ в \eqref{eq:nonlinear_ode}:

\[
  \vect{f}(\bar{\vect{x}}, \bar{\vect{u}}) =
  (v_0\cos 0,\, v_0\sin 0,\, 0,\, (v_0 - v_0)/\tau_v,\, (0-0)/\tau_\omega)\T
  = (v_0, 0, 0, 0, 0)\T.
\]

Координата $p_x$ растёт со скоростью $v_0$ --- логично; всё остальное
не меняется. Опорная точка корректна: при $\vect{x} = \bar{\vect{x}}$,
$\vect{u} = \bar{\vect{u}}$ робот действительно равномерно едет
вперёд.

\section{Ряд Тейлора}

Пусть $\Delta\vect{x} = \vect{x} - \bar{\vect{x}}$,
$\Delta\vect{u} = \vect{u} - \bar{\vect{u}}$ --- \emph{малые}
отклонения от опорной. Разложим $\vect{f}$ в ряд:

\[
  \vect{f}(\vect{x}, \vect{u}) \;=\;
    \underbrace{\vect{f}(\bar{\vect{x}}, \bar{\vect{u}})}_{\dot{\bar{\vect{x}}}}
    +
    \underbrace{\left.\frac{\partial \vect{f}}{\partial \vect{x}}\right|_*}_{\mat{A}}
    \Delta\vect{x}
    +
    \underbrace{\left.\frac{\partial \vect{f}}{\partial \vect{u}}\right|_*}_{\mat{B}}
    \Delta\vect{u}
    + O(\|\Delta\|^2).
\]

Вычитая $\dot{\bar{\vect{x}}}$ из обеих частей и обозначив
$\Delta\dot{\vect{x}} = \dot{\vect{x}} - \dot{\bar{\vect{x}}}$, отбрасывая
квадратичные члены, получаем:

\begin{formulabox}[title={Линеаризованная модель}]
\begin{equation}
  \Delta\dot{\vect{x}}(t) \;=\; \mat{A}\,\Delta\vect{x}(t)
                                 + \mat{B}\,\Delta\vect{u}(t).
  \label{eq:linear_ode}
\end{equation}
\end{formulabox}

\textbf{Важно:} формула \eqref{eq:linear_ode} верна \emph{только в окрестности}
$\bar{\vect{x}}, \bar{\vect{u}}$. Если робот сильно отклонится от опорной
траектории (например, разверулся на $90^\circ$), линейная модель станет
заметно врать.

\subsection{Вычисление матрицы $\mat{A}$}

Матрица Якоби $\mat{A} = \partial \vect{f}/\partial \vect{x}$ --- это
матрица $5 \times 5$ со столбцами $\partial \vect{f}/\partial x_j$.
Считаем частные производные функций $f_1, \ldots, f_5$ по $x_j \in \{p_x,
p_y, \theta, v, \omega\}$ и подставляем опорные значения
$\bar{\theta} = 0$, $\bar{v} = v_0$.

\textbf{Строка 1:} $f_1 = v\cos\theta$.
\[
  \frac{\partial f_1}{\partial p_x} = 0,
  \quad
  \frac{\partial f_1}{\partial p_y} = 0,
  \quad
  \frac{\partial f_1}{\partial \theta}\bigg|_* = -v\sin\theta\big|_0 = 0,
  \quad
  \frac{\partial f_1}{\partial v}\bigg|_* = \cos\theta\big|_0 = 1,
  \quad
  \frac{\partial f_1}{\partial \omega} = 0.
\]

\textbf{Строка 2:} $f_2 = v\sin\theta$.
\[
  \frac{\partial f_2}{\partial \theta}\bigg|_* = v\cos\theta\big|_0 = v_0,
  \quad
  \frac{\partial f_2}{\partial v}\bigg|_* = \sin\theta\big|_0 = 0.
\]
Остальные --- нули.

\textbf{Строка 3:} $f_3 = \omega$.
$\partial f_3/\partial \omega = 1$. Остальные --- нули.

\textbf{Строка 4:} $f_4 = (u_v - v)/\tau_v$.
$\partial f_4/\partial v = -1/\tau_v$. Остальные --- нули.

\textbf{Строка 5:} $f_5 = (u_\omega - \omega)/\tau_\omega$.
$\partial f_5/\partial \omega = -1/\tau_\omega$. Остальные --- нули.

Собираем:
\begin{formulabox}[title={Матрица $\mat{A}$ (непрерывная)}]
\begin{equation}
  \mat{A} = \begin{pmatrix}
    0 & 0 & 0   & 1         & 0          \\
    0 & 0 & v_0 & 0         & 0          \\
    0 & 0 & 0   & 0         & 1          \\
    0 & 0 & 0   & -1/\tau_v & 0          \\
    0 & 0 & 0   & 0         & -1/\tau_\omega
  \end{pmatrix}.
  \label{eq:A_continuous}
\end{equation}
\end{formulabox}

\subsection{Вычисление матрицы $\mat{B}$}

Аналогично для $\mat{B} = \partial \vect{f}/\partial \vect{u}$, размер $5 \times 2$:

\[
  \frac{\partial f_4}{\partial u_v} = \frac{1}{\tau_v},
  \quad
  \frac{\partial f_5}{\partial u_\omega} = \frac{1}{\tau_\omega}.
\]
Остальные элементы --- нули.

\begin{formulabox}[title={Матрица $\mat{B}$ (непрерывная)}]
\begin{equation}
  \mat{B} = \begin{pmatrix}
    0          & 0              \\
    0          & 0              \\
    0          & 0              \\
    1/\tau_v   & 0              \\
    0          & 1/\tau_\omega
  \end{pmatrix}.
  \label{eq:B_continuous}
\end{equation}
\end{formulabox}

Если сравнить эти матрицы с кодом \filepath{state\_space\_model.py:29}
(глава~\ref{ch:dynamics}, листинг) --- буквально одно и то же.
\emph{Это и есть наша модель.}

\subsection{Численный пример}

При $v_0 = 0.2$~м/с, $\tau_v = 0.15$~с, $\tau_\omega = 0.10$~с:

\begin{equation}
  \mat{A} = \begin{pmatrix}
    0 & 0 & 0    & 1      & 0 \\
    0 & 0 & 0.2  & 0      & 0 \\
    0 & 0 & 0    & 0      & 1 \\
    0 & 0 & 0    & -6.667 & 0 \\
    0 & 0 & 0    & 0      & -10
  \end{pmatrix},
  \quad
  \mat{B} = \begin{pmatrix}
    0      & 0 \\
    0      & 0 \\
    0      & 0 \\
    6.667  & 0 \\
    0      & 10
  \end{pmatrix}.
  \label{eq:AB_numeric}
\end{equation}

\section{Управляемость}

Перед тем как строить регулятор, надо убедиться, что робот \emph{в принципе}
может быть приведён из любого состояния в любое другое. Это свойство
\textbf{управляемости}.

\begin{formulabox}[title={Критерий Калмана}]
Объект полностью управляем тогда и только тогда, когда матрица
управляемости имеет полный ранг:
\begin{equation}
  \mat{S}_y = \big[\,\mat{B}\,\big|\,\mat{A}\mat{B}\,\big|\,\mat{A}^2\mat{B}
              \,\big|\,\ldots\,\big|\,\mat{A}^{n-1}\mat{B}\,\big],
  \qquad
  \rank \mat{S}_y = n.
  \label{eq:controllability}
\end{equation}
\end{formulabox}

Для нашей $n = 5$, $r = 2$: $\mat{S}_y$ имеет размер $5 \times 10$.
Численная проверка в MATLAB (\filepath{matlab/analyze\_system.m}) даёт
$\rank \mat{S}_y = 5$ при $v_0 \ne 0$ --- система полностью управляема.

\begin{warnbox}
\textbf{При $v_0 = 0$} опорная точка вырождена. В \eqref{eq:A_continuous}
второй столбец, отвечающий за связь $\dot{p}_y$ с $\theta$, обнуляется ---
координата $p_y$ становится неуправляемой через малые отклонения
$\Delta\theta$. Физически: стоя на месте, робот не умеет «поехать вбок»,
не повернувшись сильно (а это уже не «малые отклонения», и линейная
модель не работает).

\textbf{Следствие:} линеаризуем только вокруг $v_0 > 0$. Чтобы стартовать
с нуля, мы используем \emph{два разных режима}: TURN на стоящем роботе
управляется тем же MPC, но \emph{с другой опорной} (об этом --- глава
\ref{ch:turn_phase}).
\end{warnbox}

\section{Дискретизация: ZOH}

MPC работает с дискретной моделью, потому что:

\begin{itemize}
  \item Цикл управления --- дискретный (тик $T_s$).
  \item Команды моторам тоже подаются дискретно: «на интервале $[k T_s,
        (k+1) T_s)$ команда $\vect{u}_k$, потом следующая».
\end{itemize}

\begin{ideabox}
\textbf{ZOH (zero-order hold).} «Нулевого порядка удержание»: команда
$\vect{u}_k$ \emph{постоянна} на всём интервале $[k T_s, (k+1)T_s)$.
Это точно соответствует тому, как мы кормим $\texttt{cmd\_vel}$
моторам --- между тиками значение не меняется.

В этом предположении уравнение $\dot{\vect{x}} = \mat{A}\vect{x} +
\mat{B}\vect{u}$ можно проинтегрировать \emph{точно} от $kT_s$ до
$(k+1)T_s$ --- получим формулу с матричной экспонентой.
\end{ideabox}

\subsection{Аналитический вывод}

Решение линейной ОДУ с постоянной $\vect{u}_k$ на $[t_0, t_0 + T_s]$:

\[
  \vect{x}(t_0 + T_s)
    = e^{\mat{A} T_s}\,\vect{x}(t_0)
    + \int_0^{T_s} e^{\mat{A}(T_s - \tau)}\mat{B}\,d\tau \cdot \vect{u}_k.
\]

Сменой переменной $\sigma = T_s - \tau$ интеграл превращается в

\[
  \int_0^{T_s} e^{\mat{A}\sigma}\,d\sigma \cdot \mat{B}.
\]

Поэтому:

\begin{formulabox}[title={ZOH-дискретизация}]
\begin{equation}
  \boxed{\;
    \mat{A}_d = e^{\mat{A} T_s},
    \qquad
    \mat{B}_d = \left(\int_0^{T_s} e^{\mat{A}\sigma}\,d\sigma\right) \mat{B}.
    \;}
  \label{eq:zoh}
\end{equation}
\end{formulabox}

\subsection{Блочная экспонента (как считают в коде)}

Численно интеграл считать неудобно --- есть \textbf{трюк с блочной
матричной экспонентой}. Соберём блочную матрицу

\[
  \mat{M} =
  \begin{pmatrix}
    \mat{A} & \mat{B} \\
    \mat{0} & \mat{0}
  \end{pmatrix} \in \Real^{(n+r) \times (n+r)}.
\]

Аккуратное матричное возведение в экспоненту даёт:

\begin{equation}
  \exp(\mat{M} T_s) =
  \begin{pmatrix}
    e^{\mat{A} T_s} & \int_0^{T_s} e^{\mat{A}\sigma}d\sigma \cdot \mat{B} \\
    \mat{0}         & \mat{I}_r
  \end{pmatrix}
  =
  \begin{pmatrix}
    \mat{A}_d & \mat{B}_d \\
    \mat{0}   & \mat{I}_r
  \end{pmatrix}.
  \label{eq:block_exp}
\end{equation}

То есть: \textbf{одна} операция \texttt{expm} даёт сразу обе матрицы.

\begin{codebox}[title={\filepath{pi\_nodes/control/state\_space\_model.py:48}}]
\begin{lstlisting}[language=Python]
def zoh_discretize(A, B, Ts):
    """Exact ZOH discretization via block matrix exponential."""
    n, r = A.shape[0], B.shape[1]
    M = np.zeros((n + r, n + r))
    M[:n, :n] = A
    M[:n, n:] = B
    M_d = expm(M * Ts)
    Ad = M_d[:n, :n]
    Bd = M_d[:n, n:]
    return Ad, Bd
\end{lstlisting}
\end{codebox}

\subsection{Полученная дискретная модель}

\begin{formulabox}[title={Дискретная модель робота}]
\begin{equation}
  \boxed{\;\vect{x}_{k+1} = \mat{A}_d\,\vect{x}_k + \mat{B}_d\,\vect{u}_k.\;}
  \label{eq:discrete_model}
\end{equation}
\end{formulabox}

\textbf{Это финальная форма}, с которой будут работать ЛКР, MPC, симулятор.

\section{Собственные значения}

Собственные значения говорят, как себя ведёт открытая (без управления)
система.

\begin{notebox}
Связь между непрерывным $\lambda$ и дискретным $\tilde{\lambda}$:
\[
  \tilde{\lambda}_i = e^{\lambda_i T_s}.
\]
Устойчивость:
\begin{itemize}
  \item непрерывная: $\operatorname{Re}\lambda_i < 0$ (левая полуплоскость);
  \item дискретная: $|\tilde{\lambda}_i| < 1$ (внутри единичного круга).
\end{itemize}
\end{notebox}

Для нашей $\mat{A}$ \eqref{eq:A_continuous} собственные значения легко
считать «руками» --- $\mat{A}$ блочно-треугольная:

\[
  \det(\mat{A} - \lambda \mat{I}) =
  (-\lambda)^3 \cdot (-1/\tau_v - \lambda) \cdot (-1/\tau_\omega - \lambda).
\]

\begin{equation}
  \lambda_{1,2,3} = 0,
  \qquad
  \lambda_4 = -1/\tau_v,
  \qquad
  \lambda_5 = -1/\tau_\omega.
  \label{eq:eigenvalues}
\end{equation}

Три нуля --- это \emph{интеграторы}: $p_x, p_y, \theta$ накапливаются
во времени без внутреннего «возврата к нулю». Два отрицательных корня
--- инерционность моторов: возмущение по скорости затухает.

В дискретной форме: $\tilde\lambda_{1,2,3} = 1$ (на границе устойчивости),
$\tilde\lambda_4 = e^{-T_s/\tau_v} \approx 0.717$,
$\tilde\lambda_5 = e^{-T_s/\tau_\omega} \approx 0.607$.

\begin{warnbox}
\textbf{Объект не асимптотически устойчив.} Три собственных значения
на единичной окружности означают: без управления отклонения по координатам
и курсу не затухают (что логично --- робот не «возвращается домой» сам
по себе). Это норма для движущихся объектов. ЛКР/MPC сделают замкнутую
систему устойчивой.
\end{warnbox}

\section{Каноническая МПС-форма $(s, v, \theta, \omega, e_\text{int})$}

Выше мы работали с легаси-формой $(p_x, p_y, \theta, v, \omega)$,
потому что её удобно линеаризовать «с нуля» из физики. В коде
\filepath{mps\_node.py} используется \emph{другая}, каноническая форма
$(s, v, \theta, \omega, e_\text{int})$. Концептуально это та же модель,
просто:

\begin{itemize}
  \item вместо двух декартовых $(p_x, p_y)$ --- одна продольная $s$
        (мы заранее знаем, что в сценарии $p_y \equiv 0$);
  \item плюс интегральный член $e_\text{int}$, чтобы устранить
        установившуюся ошибку скорости.
\end{itemize}

Конкретные матрицы $\mat{A}, \mat{B}$ для этой формы синтезируются
в \filepath{matlab/main.m} и экспортируются в \filepath{config.yaml},
секция \texttt{mps.matrices.\{A,B,C,D\}}. Размер тот же ---  $5 \times 5$ и $5
\times 2$, но числа другие. Период тоже \emph{другой}: $T_s = 0.02$~с
(50~Гц), а не $0.05$~с легаси-контура.

\begin{notebox}
В коде эти матрицы приходят в \emph{непрерывной} форме (контракт
\filepath{docs/mps/api.md}). \texttt{mps\_node} ZOH-дискретизирует их
при загрузке/Apply --- так что одна и та же $\mat{A}$ работает
и в симе на ноуте, и на Pi, в любых обновлениях.
\end{notebox}

\section{Что мы получили}

\begin{itemize}
  \item \textbf{Линеаризация}: ряд Тейлора 1-го порядка вокруг опорного
        $(v_0, 0)\T$. Получились $\mat{A} \in \Real^{5\times 5}$,
        $\mat{B} \in \Real^{5\times 2}$ с явными формулами \eqref{eq:A_continuous},
        \eqref{eq:B_continuous}.
  \item Объект \textbf{полностью управляем} при $v_0 \ne 0$.
  \item \textbf{Дискретизация ZOH}: $\mat{A}_d = e^{\mat{A}T_s}$, $\mat{B}_d$
        --- через интеграл. Вычисляется одним вызовом \texttt{expm} от
        блочной матрицы.
  \item Собственные значения непрерывной системы: три нуля (интеграторы
        по $p_x, p_y, \theta$) и два отрицательных. Открытый объект не
        асимптотически устойчив, но управляем.
  \item Каноническая МПС-форма --- та же модель с $(s, \cdot, \theta, \omega,
        e_\text{int})$ вместо $(p_x, p_y, \theta, v, \omega)$ и
        $T_s = 0.02$~с.
\end{itemize}

\medskip
\textbf{Теперь у нас есть всё, что нужно для регулятора.} Дальше начинаем
строить управляющий контур.

% ============================================================
%  Глава 6 — Что такое регулятор: от P к ЛКР
% ============================================================
\chapter{Что такое регулятор}
\label{ch:controllers}

\section{Вопрос}

У нас есть модель $\vect{x}_{k+1} = \mat{A}_d\vect{x}_k + \mat{B}_d\vect{u}_k$.
Робот «понимает», как состояние реагирует на команду. \textbf{Какую
команду $\vect{u}_k$ ему подавать на каждом тике}, чтобы достичь цели?

Эта задача --- задача синтеза \emph{регулятора}.

\begin{ideabox}
Регулятор --- это \emph{функция}, которая по текущему состоянию (и,
возможно, желаемому состоянию) выдаёт управление:
\[
  \vect{u}_k = \pi(\vect{x}_k, \vect{x}^\text{ref}_k).
\]
Слово «регулятор» --- из ТАУ. Часто говорят также «контроллер»,
«controller», «policy». Это одно и то же.
\end{ideabox}

Идея этой главы --- \emph{интуитивно} объяснить, чем отличаются простой
П-регулятор, PID и ЛКР. К концу главы будет понятно, \emph{зачем} мы вообще
идём в матрицы $\mat{Q}, \mat{R}$ и почему нельзя обойтись пропорциональным
коэффициентом.

\section{П-регулятор (proportional)}

Самый простой регулятор: чем дальше от цели --- тем сильнее тяни.

\begin{equation}
  u_k = -k_p \cdot (x_k - x^\text{ref}_k).
  \label{eq:P_controller}
\end{equation}

Здесь $k_p > 0$ --- \emph{коэффициент усиления} (gain). Чем он больше,
тем «жёстче» регулятор. Скаляр; в многомерном случае --- вектор/матрица.

\begin{examplebox}[title={Скаляр, удержание скорости}]
Хотим, чтобы робот ехал ровно $v_\text{target}$. Текущая скорость $v$.
Тогда $u = -k_p (v - v_\text{target})$. Если $v < v_\text{target}$ ---
команда положительная, разгон. Если $v > v_\text{target}$ --- отрицательная,
торможение.
\end{examplebox}

\textbf{Проблемы П-регулятора:}

\begin{itemize}
  \item \textbf{Установившаяся ошибка.} Если есть сопротивление (трение,
        уклон), при достижении $v = v_\text{target}$ команда обнуляется,
        но трение остаётся --- робот замедлится. На равновесии останется
        ошибка $\ne 0$.
  \item \textbf{Колебания.} Большой $k_p$ $\to$ перерегулирование $\to$
        качается вокруг цели.
  \item \textbf{Не учитывает динамику.} П-регулятор «слепой»: ему всё равно,
        что моторы инерционны. Если $\tau_v$ велико --- регулятор
        будет «опаздывать» на каждом тике.
\end{itemize}

\section{PID}

Классическая попытка починить проблемы П: добавить \emph{интеграл} ошибки
(чтобы убрать установившуюся) и \emph{производную} (чтобы предугадать).

\begin{equation}
  u_k = -k_p \cdot e_k - k_i \cdot \sum_{j=0}^{k} e_j \Delta t - k_d \cdot (e_k - e_{k-1})/\Delta t.
  \label{eq:PID_controller}
\end{equation}

PID --- старший товарищ. Работает на всём: от лифта до автопилота. Но
у него тоже минусы:

\begin{itemize}
  \item Подбор $k_p, k_i, k_d$ --- искусство (Зиглер--Никольс, ручное
        крутение). На многомерных задачах перестаёт быть удобным.
  \item Не оптимален по никакому критерию: «работает» $\ne$ «работает
        лучше всего».
  \item \textbf{Не знает модель} объекта. Если бы знал $\mat{A}_d, \mat{B}_d$
        --- мог бы предсказывать, а не реагировать. PID не предсказывает.
\end{itemize}

\section{Линейно-квадратичный регулятор (ЛКР)}

ЛКР (linear-quadratic regulator, LQR) --- иначе подходит к задаче.

\begin{ideabox}
\textbf{Идея ЛКР:} вместо ручного подбора коэффициентов
\emph{сформулировать} цель как минимизацию специальной функции
\textbf{критерия качества}, и пусть математика сама найдёт лучший
регулятор для этой цели.
\end{ideabox}

Критерий --- сумма «штрафов» по всему будущему:

\begin{formulabox}[title={ЛКР-критерий качества}]
\begin{equation}
  J = \sum_{k=0}^{\infty} \Big(\vect{x}_k\T \mat{Q}\,\vect{x}_k +
                                  \vect{u}_k\T \mat{R}\,\vect{u}_k\Big).
  \label{eq:lqr_cost}
\end{equation}
\end{formulabox}

Здесь:

\begin{itemize}
  \item $\mat{Q} \in \Real^{n\times n}$ --- штраф за \textbf{отклонения
        состояния} от нуля (предполагаем, что цель --- $\vect{x} = 0$;
        если цель ненулевая, перейдём в координаты ошибки).
  \item $\mat{R} \in \Real^{r\times r}$ --- штраф за \textbf{величину
        управления}. Большой $\mat{R}$ --- регулятор будет «экономить»
        моторы, малый --- «топить педаль в пол».
  \item Обе матрицы --- симметричные, положительно (полу-)определённые.
        На практике диагональные:
        $\mat{Q} = \diag(q_1, \ldots, q_n)$, $\mat{R} = \diag(r_1, \ldots, r_r)$.
\end{itemize}

\subsection{Что значит «положительно определённая»}

\begin{notebox}
Матрица $\mat{Q}$ \textbf{положительно (полу-)определённая}, если для
любого ненулевого $\vect{x}$ выполняется $\vect{x}\T \mat{Q}\,\vect{x} > 0$
(или $\ge 0$ --- полу-).

Для диагональной $\mat{Q} = \diag(q_1, \ldots, q_n)$ это сводится к
$q_i \ge 0$ (или строгая положительность для всех $q_i$, тогда уже
positive definite).

Зачем требование: штраф $\vect{x}\T \mat{Q}\,\vect{x}$ должен быть
\emph{неотрицательным}, иначе оптимизатор «разгонит» состояние до
$-\infty$ и получит бесконечную выгоду. Аналогично $\mat{R}$ должна
быть строго положительно определённая --- иначе можно бесконечно
увеличивать $\vect{u}$ без штрафа.
\end{notebox}

\subsection{Зачем квадратичный штраф?}

Почему именно $\vect{x}\T\mat{Q}\vect{x}$, а не, скажем, $|\vect{x}|$?

\begin{itemize}
  \item \textbf{Дифференцируем} $\to$ можно решить аналитически.
  \item Маленькие ошибки штрафуются \emph{ещё меньше} (квадрат малого числа
        --- ещё меньше), большие --- \emph{гораздо больше} (квадрат
        $5 = 25$, квадрат $10 = 100$). Регулятор сильно реагирует на
        большие отклонения, мягко --- на мелкие.
  \item Удобная аналогия с энергией / расстоянием. Норма $\|\vect{x}\|^2
        = \vect{x}\T\vect{x}$, а $\vect{x}\T\mat{Q}\vect{x}$ --- «взвешенная»
        норма.
\end{itemize}

\subsection{Результат ЛКР}

\textbf{Ключевой факт.} Если объект линейный, дискретный, управляемый,
и критерий \eqref{eq:lqr_cost} с положительно определёнными $\mat{Q}, \mat{R}$,
то оптимальный регулятор --- \textbf{линейная обратная связь}:

\begin{formulabox}[title={ЛКР-закон управления}]
\begin{equation}
  \vect{u}_k = -\mat{K}\,\vect{x}_k,
  \label{eq:lqr_law}
\end{equation}
\end{formulabox}

где $\mat{K}$ --- матрица $r \times n$, не зависящая от $k$ (постоянная).
Она вычисляется \emph{один раз}, оффлайн, через \emph{уравнение Риккати}.
\emph{Как именно} --- следующая глава.

\section{Интуитивная карта}

\begin{center}
\begin{tabular}{lp{8cm}}
\toprule
\textbf{Регулятор} & \textbf{Идея}\\
\midrule
П    & «Тяни сильнее, если дальше».\\
PID  & П + интеграл + производная (борьба с трением и инерцией).\\
ЛКР  & Минимизируй квадратичный критерий $J = \sum (\vect{x}\T \mat{Q}\vect{x}
       + \vect{u}\T \mat{R}\vect{u})$. Решение --- $\vect{u} = -\mat{K}\vect{x}$.
       \textbf{Оптимален по выбранному критерию.}\\
MPC  & Каждый тик решай ЛКР-подобную задачу \emph{на короткий горизонт
       вперёд}, с \emph{учётом ограничений} $\vect{u}_\text{min} \le \vect{u} \le \vect{u}_\text{max}$.\\
\bottomrule
\end{tabular}
\end{center}

\begin{ideabox}
\textbf{Логическая цепочка:}
\begin{itemize}
  \item П --- слепая реакция.
  \item PID --- реакция с памятью и предсказанием в эвристике.
  \item ЛКР --- \emph{оптимальная} реакция при заданном критерии и
        \emph{знании модели}. Не учитывает ограничения.
  \item MPC --- ЛКР с учётом ограничений на $\vect{u}$ и предсказанием
        на конечный горизонт $N$.
\end{itemize}
Каждый следующий --- более мощный и при этом более вычислительно дорогой.
В проекте используются ЛКР и MPC одновременно (MPC при \texttt{solver=clip}
сводится к ЛКР с клиппингом --- эффективный компромисс).
\end{ideabox}

\section{Связь с проектом}

\begin{codebox}[title={В проекте есть оба:}]
\begin{itemize}
  \item \filepath{pi\_nodes/control/lqr\_controller.py} --- чистый ЛКР
        (используется как baseline и для сравнения).
  \item \filepath{pi\_nodes/control/mpc\_controller.py} --- MPC. В режиме
        \texttt{solver = "clip"} он \emph{буквально} вычисляет ЛКР-усиление
        $\mat{K}_\text{first}$ (первая строка лифтированной QP-задачи),
        применяет его и клиппит в ограничения.
\end{itemize}
\end{codebox}

В следующих двух главах разберём ЛКР, потом MPC.

% ============================================================
%  Глава 7 — ЛКР через DARE
% ============================================================
\chapter{ЛКР: оптимальный регулятор}
\label{ch:lqr}

\section{Вопрос}

В прошлой главе мы сказали: ЛКР даёт $\vect{u}_k = -\mat{K}\vect{x}_k$,
и $\mat{K}$ считается через \emph{уравнение Риккати}. Сейчас аккуратно
разберём, \emph{что} это за уравнение и \emph{откуда} оно берётся.

\begin{ideabox}
ЛКР --- это \emph{задача оптимизации} на бесконечном горизонте.
Решить её «в лоб» (например, методом Беллмана) можно, но получается
тяжёлая рекурсия. Чудо в том, что \textbf{для линейного объекта с
квадратичным критерием эта рекурсия сходится к одной матрице $\mat{P}$,
которую можно найти из алгебраического уравнения}.

Это уравнение --- DARE (Discrete-time Algebraic Riccati Equation).
\end{ideabox}

\section{Постановка задачи}

Объект:
\[
  \vect{x}_{k+1} = \mat{A}_d\,\vect{x}_k + \mat{B}_d\,\vect{u}_k,
  \qquad \vect{x}_0 \text{ задано}.
\]

Критерий:
\[
  J = \sum_{k=0}^{\infty} \Big(\vect{x}_k\T \mat{Q}\,\vect{x}_k +
                                  \vect{u}_k\T \mat{R}\,\vect{u}_k\Big).
\]

Цель: найти последовательность $\{\vect{u}_k\}_{k=0}^\infty$, которая
\textbf{минимизирует} $J$.

\section{Принцип динамического программирования}

\begin{ideabox}
\textbf{Принцип Беллмана.} Оптимальная политика на бесконечном
горизонте обладает свойством: \emph{какова бы ни была начальная
ситуация и первое решение, последующие решения должны быть
оптимальными для оставшейся задачи}.

Для квадратичного критерия и линейной динамики из этого принципа
следует, что \textbf{функция ценности} (стоимость до бесконечности
из состояния $\vect{x}$) тоже квадратичная:
\[
  V^*(\vect{x}) = \vect{x}\T \mat{P}\,\vect{x},
\]
где $\mat{P}$ --- симметричная положительно определённая матрица.
\end{ideabox}

Подставляем в уравнение Беллмана:

\[
  \vect{x}\T \mat{P}\,\vect{x}
  = \min_{\vect{u}}
    \Big\{\vect{x}\T \mat{Q}\,\vect{x} + \vect{u}\T \mat{R}\,\vect{u}
          + (\mat{A}_d\vect{x} + \mat{B}_d\vect{u})\T \mat{P}\,
            (\mat{A}_d\vect{x} + \mat{B}_d\vect{u})\Big\}.
\]

Минимум по $\vect{u}$ ищется аналитически: производная по $\vect{u}$,
равная нулю.

\subsection{Шаг 1: найти оптимальное \texorpdfstring{$\vect{u}^*$}{u*}}

\begin{align*}
  \frac{\partial}{\partial \vect{u}}\Big[
    \vect{u}\T \mat{R}\vect{u} + (\mat{A}_d\vect{x} + \mat{B}_d\vect{u})\T
    \mat{P}(\mat{A}_d\vect{x} + \mat{B}_d\vect{u})
  \Big]
  &= 2\mat{R}\vect{u} + 2\mat{B}_d\T \mat{P}\,(\mat{A}_d\vect{x} + \mat{B}_d\vect{u})\\
  &= 0.
\end{align*}

Раскрываем:
\[
  (\mat{R} + \mat{B}_d\T \mat{P}\mat{B}_d)\,\vect{u} =
    -\mat{B}_d\T \mat{P}\mat{A}_d\,\vect{x}.
\]

\begin{formulabox}[title={Оптимальное управление через $\mat{P}$}]
\begin{equation}
  \vect{u}^* = -\mat{K}\,\vect{x},
  \qquad
  \mat{K} = (\mat{R} + \mat{B}_d\T \mat{P}\mat{B}_d)^{-1} \mat{B}_d\T \mat{P}\mat{A}_d.
  \label{eq:lqr_gain}
\end{equation}
\end{formulabox}

Это \textbf{линейный} закон: оптимальное управление --- линейная функция
текущего состояния. Чудо ЛКР именно в этом.

\subsection{Шаг 2: уравнение для $\mat{P}$ (DARE)}

Подставим $\vect{u}^* = -\mat{K}\vect{x}$ обратно в уравнение Беллмана
и упростим. После выкладок получим:

\begin{formulabox}[title={DARE --- Discrete Algebraic Riccati Equation}]
\begin{equation}
  \mat{P} \;=\; \mat{A}_d\T \mat{P}\mat{A}_d
              - \mat{A}_d\T \mat{P}\mat{B}_d
                (\mat{R} + \mat{B}_d\T \mat{P}\mat{B}_d)^{-1}
                \mat{B}_d\T \mat{P}\mat{A}_d
              + \mat{Q}.
  \label{eq:dare}
\end{equation}
\end{formulabox}

Это \emph{нелинейное матричное уравнение} относительно $\mat{P}$.

\section{Как решают DARE численно}

\begin{notebox}
Аналитически решить DARE невозможно (кроме скалярных случаев). На практике
используют численные методы:

\begin{itemize}
  \item \textbf{Итерации Риккати} --- запускаем рекурсию
        $\mat{P}_{k+1} = \text{RHS DARE}(\mat{P}_k)$ от любой начальной
        $\mat{P}_0$ (часто $\mat{Q}$). Сходится к решению, если объект
        управляем и $\mat{Q}, \mat{R}$ положительны.
  \item \textbf{Метод Шура} (Шурова декомпозиция гамильтоновой матрицы) ---
        стандарт в MATLAB и SciPy.
\end{itemize}

В нашем проекте --- наивные итерации (см. \filepath{pi\_nodes/control/\_linalg.py}),
fallback на \texttt{scipy.linalg.solve\_discrete\_are}, когда scipy есть.
\end{notebox}

\begin{codebox}[title={\filepath{pi\_nodes/control/\_linalg.py}, функция \texttt{solve\_discrete\_are}}]
\begin{lstlisting}[language=Python]
def solve_discrete_are(A, B, Q, R, max_iter=10_000, tol=1e-9):
    """Решает DARE итерациями Беллмана."""
    P = Q.copy()
    for _ in range(max_iter):
        BtP = B.T @ P
        K = np.linalg.solve(R + BtP @ B, BtP @ A)
        P_new = A.T @ P @ A - A.T @ P @ B @ K + Q
        if np.max(np.abs(P_new - P)) < tol:
            return P_new
        P = P_new
    raise RuntimeError("DARE did not converge")
\end{lstlisting}
\end{codebox}

После того как DARE решено и $\mat{P}$ найдена, считаем
$\mat{K}$ по \eqref{eq:lqr_gain}. Готово --- регулятор готов к работе.

\section{Какие $\mat{Q}, \mat{R}$ выбрать}

\begin{ideabox}
Выбор $\mat{Q}, \mat{R}$ --- это \emph{инженерное искусство}. ЛКР не
говорит, какие они должны быть; он только обещает \emph{оптимальность}
для тех, которые мы зададим.
\end{ideabox}

Практические рецепты:

\subsection{Правило Брайсона}

Для каждой компоненты $i$ положить $q_i = 1/x_{i,\max}^2$, где
$x_{i,\max}$ --- максимально \emph{допустимое} отклонение. Аналогично
для $r_j = 1/u_{j,\max}^2$. Это нормализует масштабы --- регулятор не
будет «думать», что отклонение в радианах ($\sim 0.1$) важнее, чем в
метрах ($\sim 5$).

\subsection{Дефолты проекта}

\begin{codebox}[title={\filepath{config.yaml}, секция \texttt{control.weights}}]
\begin{verbatim}
control:
  weights:
    Q_diag: [10.0, 10.0, 5.0, 1.0, 1.0]
    R_diag: [1.0, 1.0]
\end{verbatim}
\end{codebox}

\begin{tabular}{lll}
\toprule
Вес & Компонента & Смысл\\
\midrule
$10$ & $p_x$       & сильный штраф за «недокат» / «перекат»\\
$10$ & $p_y$       & сильный штраф за «съезд вбок»\\
$5$  & $\theta$    & средний штраф за кривой курс\\
$1$  & $v$         & лёгкий штраф за «не та скорость»\\
$1$  & $\omega$    & лёгкий штраф за «не то вращение»\\
$1$  & $u_v$       & вес на расход линейного управления\\
$1$  & $u_\omega$  & вес на расход углового управления\\
\bottomrule
\end{tabular}

\subsection{Что меняется, если\dots}

\begin{whatifbox}
\textbf{Q $\gg$ R} --- регулятор готов «топить педаль в пол», лишь бы
$\vect{x}$ был у цели. Быстро, но энергоёмко и с перерегулированием.

\textbf{Q $\ll$ R} --- регулятор «бережёт моторы», медленно подходит
к цели. Гладко, но долго.

\textbf{$\mat{Q}$ не-диагональная} --- разрешает «обмен» между
компонентами: можно сильно штрафовать, например, разницу $p_x - s_\text{ref}$,
а не каждую компоненту отдельно. Усложняет интерпретацию, редко нужно.

\textbf{$q_\theta = 0$} --- регулятору всё равно, куда смотрит робот.
Странно, но математически валидно: ЛКР решит DARE и выдаст $\mat{K}$,
где курс будет «болтаться» свободно.

\textbf{$r_i = 0$} --- неуправляемо, DARE не сходится, $\mat{P}$
вырождается. \emph{Не делать так.}
\end{whatifbox}

\section{Замкнутая система}

\begin{ideabox}
Подставим закон $\vect{u}_k = -\mat{K}\vect{x}_k$ в исходную модель:
\[
  \vect{x}_{k+1} = \mat{A}_d\vect{x}_k + \mat{B}_d(-\mat{K}\vect{x}_k)
                 = (\mat{A}_d - \mat{B}_d\mat{K})\,\vect{x}_k.
\]
Это \textbf{замкнутая система}. Её собственные значения $\lambda(\mat{A}_d
- \mat{B}_d \mat{K})$ --- \emph{полюса} замкнутого контура. ЛКР гарантирует,
что все они лежат внутри единичного круга ($|\lambda| < 1$) --- замкнутая
система \emph{устойчива}.
\end{ideabox}

Проверка в коде:

\begin{codebox}[title={\filepath{pi\_nodes/control/mpc\_controller.py:332}}]
\begin{lstlisting}[language=Python]
def closed_loop_eigenvalues(self) -> np.ndarray:
    """|λ| < 1 ⇔ closed-loop x_{k+1} = (Ad - Bd K) x is stable."""
    return np.linalg.eigvals(self.Ad - self.Bd @ self.K_first)

def is_stable(self) -> bool:
    return bool(np.all(np.abs(self.closed_loop_eigenvalues()) < 1.0 - 1e-9))
\end{lstlisting}
\end{codebox}

\section{Слабое место: ограничения}

ЛКР \textbf{не знает} о том, что $u_v \le 0.30$~м/с и $|u_\omega| \le
2$~рад/с. Если задать $\vect{x}_0$ далеко от цели, закон $-\mat{K}\vect{x}_0$
выдаст $\vect{u}$, который физически невозможен.

\begin{warnbox}
В коде \texttt{LQRController.step()} это решают \emph{пост-обработкой}:
после $-\mat{K}\vect{x}$ результат прогоняется через
\texttt{np.clip(u, u\_min, u\_max)}. Это работает, но \emph{не оптимально}
--- ЛКР проектировал под полную свободу. Когда $\vect{u}$ упирается в
насыщение, его поведение перестаёт быть оптимальным.

Это --- \emph{мотивация для MPC}. MPC формулирует ограничения внутри
оптимизации, а не накладывает их сверху.
\end{warnbox}

\section{Что мы получили}

\begin{itemize}
  \item \textbf{ЛКР} --- линейный оптимальный регулятор:
        $\vect{u}_k = -\mat{K}\vect{x}_k$.
  \item Матрица $\mat{K}$ считается \emph{оффлайн} через DARE
        \eqref{eq:dare} --- одно матричное уравнение, решается итерациями
        либо Шуром.
  \item ЛКР гарантирует устойчивость замкнутой системы (все полюса
        внутри единичного круга).
  \item Выбор $\mat{Q}, \mat{R}$ --- инженерное искусство: правило
        Брайсона, диагональные веса.
  \item ЛКР \textbf{не знает} о box-ограничениях на $\vect{u}$. Клип
        сверху --- костыль, мотивация для MPC.
\end{itemize}

Идём дальше: MPC.

% ============================================================
%  Глава 8 — MPC: горизонт, лифтинг, QP
% ============================================================
\chapter{MPC: Model Predictive Control}
\label{ch:mpc}

\section{Вопрос}

ЛКР минимизирует $J$ на \emph{бесконечном} горизонте, выдаёт постоянное
усиление $\mat{K}$ и не знает о ограничениях $\vect{u}_\text{min} \le
\vect{u} \le \vect{u}_\text{max}$. Эти два недостатка устраняет MPC.

\begin{ideabox}
\textbf{Идея MPC:} вместо одного решения на $\infty$ --- решать
\emph{маленькую} оптимизационную задачу \emph{на каждом тике}, на
конечный горизонт $N$ вперёд, с явными ограничениями.

\textbf{Receding horizon (скользящий горизонт):}
\begin{enumerate}
  \item Сейчас, в момент $k$, посмотри на $N$ шагов вперёд.
  \item Сосчитай \emph{весь будущий план} $\vect{u}_0, \vect{u}_1, \ldots,
        \vect{u}_{N-1}$, минимизируя стоимость.
  \item Примени \emph{только первый шаг} $\vect{u}_0$ к роботу.
  \item На следующем тике начни заново с $\vect{x}_{k+1}$.
\end{enumerate}
Это похоже на шахматы: считаем дерево на N ходов, делаем один, считаем
заново.
\end{ideabox}

\section{Постановка}

\begin{formulabox}[title={MPC-задача (с ограничениями)}]
\begin{equation}
  \begin{aligned}
  \min_{\,\vect{u}_0, \ldots, \vect{u}_{N-1}}\;
  & J = \sum_{i=0}^{N-1}\Big(\vect{x}_i\T \mat{Q}\vect{x}_i +
                              \vect{u}_i\T \mat{R}\vect{u}_i\Big)
        + \vect{x}_N\T \mat{P}_f\,\vect{x}_N\\
  \text{при условиях } &
        \vect{x}_{i+1} = \mat{A}_d\vect{x}_i + \mat{B}_d\vect{u}_i,
        \quad i = 0, \ldots, N-1,\\
  & \vect{u}_\text{min} \le \vect{u}_i \le \vect{u}_\text{max},
        \quad i = 0, \ldots, N-1,\\
  & \vect{x}_0 \text{ задано (измерено сейчас)}.
  \end{aligned}
  \label{eq:mpc_problem}
\end{equation}
\end{formulabox}

Новое:
\begin{itemize}
  \item Сумма по \emph{конечному} $N$ (не до бесконечности).
  \item \textbf{Терминальный штраф} $\mat{P}_f$ --- штраф за последнее
        состояние $\vect{x}_N$. Если выбрать $\mat{P}_f$ как решение DARE
        для $(\mat{A}_d, \mat{B}_d, \mat{Q}, \mat{R})$, то конечный горизонт
        ведёт себя «как бесконечный» --- сумма от $N$ до $\infty$
        заменяется аккуратно квадратичным выражением.
  \item Явные box-ограничения на $\vect{u}_i$.
\end{itemize}

\section{Лифтированная динамика}

Переменными оптимизации в \eqref{eq:mpc_problem} являются $\vect{u}_0, \ldots,
\vect{u}_{N-1}$ --- всего $rN$ скаляров. Состояние $\vect{x}_i$ не
независимая переменная: оно \emph{зависит} от $\vect{u}_i$ через
$\vect{x}_{i+1} = \mat{A}_d\vect{x}_i + \mat{B}_d\vect{u}_i$. Удобно это
зависимое представление сделать явным.

\begin{ideabox}
\textbf{Лифтинг (lifting).} Разворачиваем рекурсию по динамике.
\begin{align*}
  \vect{x}_1 &= \mat{A}_d\vect{x}_0 + \mat{B}_d\vect{u}_0,\\
  \vect{x}_2 &= \mat{A}_d\vect{x}_1 + \mat{B}_d\vect{u}_1
              = \mat{A}_d^2\vect{x}_0 + \mat{A}_d\mat{B}_d\vect{u}_0 + \mat{B}_d\vect{u}_1,\\
  \vect{x}_3 &= \mat{A}_d^3\vect{x}_0 + \mat{A}_d^2\mat{B}_d\vect{u}_0
              + \mat{A}_d\mat{B}_d\vect{u}_1 + \mat{B}_d\vect{u}_2,\\
             & \vdots
\end{align*}
Складываем все состояния в большой столбец $\vect{X}$ и все управления
--- в $\vect{U}$:
\end{ideabox}

Введём блочные векторы:
\[
  \vect{X} = \begin{pmatrix} \vect{x}_1\\ \vect{x}_2\\ \vdots\\ \vect{x}_N \end{pmatrix} \in \Real^{nN},
  \qquad
  \vect{U} = \begin{pmatrix} \vect{u}_0\\ \vect{u}_1\\ \vdots\\ \vect{u}_{N-1} \end{pmatrix} \in \Real^{rN}.
\]

Тогда вся динамика «помещается» в одно матричное равенство:

\begin{formulabox}[title={Лифтированная динамика}]
\begin{equation}
  \vect{X} \;=\; \mat{\Phi}\,\vect{x}_0 \;+\; \mat{\Gamma}\,\vect{U},
  \label{eq:lifted}
\end{equation}
где
\[
  \mat{\Phi} = \begin{pmatrix}
    \mat{A}_d\\ \mat{A}_d^2\\ \vdots\\ \mat{A}_d^N
  \end{pmatrix} \in \Real^{nN \times n},
  \qquad
  \mat{\Gamma} = \begin{pmatrix}
    \mat{B}_d & 0 & \cdots & 0\\
    \mat{A}_d\mat{B}_d & \mat{B}_d & \cdots & 0\\
    \vdots & & \ddots & \vdots\\
    \mat{A}_d^{N-1}\mat{B}_d & \mat{A}_d^{N-2}\mat{B}_d & \cdots & \mat{B}_d
  \end{pmatrix} \in \Real^{nN \times rN}.
\]
\end{formulabox}

\begin{codebox}[title={\filepath{pi\_nodes/control/mpc\_controller.py:134}}]
\begin{lstlisting}[language=Python]
def _build_qp_matrices(self):
    n, r, N = self.n, self.r, self.N
    Phi = np.zeros((n * N, n))
    Gamma = np.zeros((n * N, r * N))
    Apow = np.eye(n)
    for i in range(N):
        Apow = self.Ad @ Apow                    # A^{i+1}
        Phi[i*n:(i+1)*n, :] = Apow
        for j in range(i + 1):
            block = np.linalg.matrix_power(self.Ad, i - j) @ self.Bd
            Gamma[i*n:(i+1)*n, j*r:(j+1)*r] = block
    # ... дальше Q_bar, R_bar, H, f_mat
\end{lstlisting}
\end{codebox}

\subsection{Стоимость в лифтированной форме}

Соберём блочно-диагональные веса (последний блок --- терминальный):

\[
  \bar{\mat{Q}} = \begin{pmatrix}
    \mat{Q} &       &        & \\
            & \mat{Q} &      & \\
            &       & \ddots & \\
            &       &        & \mat{P}_f
  \end{pmatrix} \in \Real^{nN\times nN},
  \qquad
  \bar{\mat{R}} = \mat{I}_N \otimes \mat{R} \in \Real^{rN\times rN}.
\]

Подставим \eqref{eq:lifted} в стоимость $J$:

\[
  J(\vect{U}) = (\mat{\Phi}\vect{x}_0 + \mat{\Gamma}\vect{U})\T \bar{\mat{Q}}
                 (\mat{\Phi}\vect{x}_0 + \mat{\Gamma}\vect{U})
               + \vect{U}\T \bar{\mat{R}}\,\vect{U}.
\]

Раскрываем и группируем по степеням $\vect{U}$:

\begin{formulabox}[title={Квадратичная форма стоимости}]
\begin{equation}
  J(\vect{U}) = \half\,\vect{U}\T \mat{H}\,\vect{U} + \vect{f}(\vect{x}_0)\T \vect{U} + \text{const},
  \label{eq:quad_cost}
\end{equation}
\[
  \mat{H} = 2\,(\mat{\Gamma}\T \bar{\mat{Q}}\mat{\Gamma} + \bar{\mat{R}}),
  \qquad
  \vect{f}(\vect{x}_0) = 2\,\mat{\Gamma}\T \bar{\mat{Q}}\mat{\Phi}\,\vect{x}_0.
\]
\end{formulabox}

«const» --- слагаемое, зависящее только от $\vect{x}_0$. Для оптимизации
по $\vect{U}$ оно неважно.

\section{Квадратичное программирование (QP)}

С учётом ограничений MPC превращается в каноническую \textbf{QP-задачу}:

\begin{formulabox}[title={QP-постановка MPC}]
\begin{equation}
  \min_{\vect{U}}\; \half \vect{U}\T \mat{H}\vect{U} + \vect{f}\T \vect{U},
  \qquad
  \vect{U}_\text{min} \le \vect{U} \le \vect{U}_\text{max}.
  \label{eq:qp}
\end{equation}
\end{formulabox}

Здесь $\vect{U}_\text{min}, \vect{U}_\text{max}$ --- повторение
$\vect{u}_\text{min}, \vect{u}_\text{max}$ по $N$ блокам.

\textbf{Свойства задачи:}
\begin{itemize}
  \item $\mat{H} = \mat{H}\T$, положительно определённая
        $\Rightarrow$ задача \textbf{выпуклая}, имеет единственный
        минимум.
  \item Ограничения --- линейные неравенства, выпуклые.
  \item Существуют десятки QP-солверов (OSQP, qpOASES, scipy.optimize и др.).
\end{itemize}

\section{Два режима в проекте: \texttt{clip} и \texttt{qp}}

В \filepath{pi\_nodes/control/mpc\_controller.py} реализованы \textbf{две}
стратегии решения:

\subsection{\texttt{clip} (по умолчанию)}

\textbf{Без} ограничений (\emph{в задаче, не в железе}) решение QP
\eqref{eq:qp} аналитическое --- производная по $\vect{U}$, равная нулю:

\[
  \mat{H}\,\vect{U}^* + \vect{f} = 0
  \;\;\Rightarrow\;\;
  \vect{U}^* = -\mat{H}^{-1}\vect{f}(\vect{x}_0).
\]

Берём первые $r$ компонент --- это $\vect{u}_0$. После лишних выкладок
получается, что

\begin{equation}
  \vect{u}_0 = -\mat{K}_\text{first}\,\vect{x}_0,
  \qquad
  \mat{K}_\text{first} = (\mat{H}^{-1}\,2\mat{\Gamma}\T \bar{\mat{Q}}\mat{\Phi})_{[1\ldots r,\,:]}.
  \label{eq:K_first}
\end{equation}

То есть аналитический MPC --- это \textbf{та же линейная обратная связь},
что ЛКР, только усиление $\mat{K}_\text{first}$ другое (зависит от $N$).
После умножения мы \emph{клиппим} управление к ограничениям:

\[
  \vect{u}_0 \leftarrow \operatorname{clip}(\vect{u}_0, \vect{u}_\text{min}, \vect{u}_\text{max}).
\]

\begin{notebox}
\textbf{Когда это оптимально:} ровно когда невозмущённое решение
$-\mat{K}_\text{first}\vect{x}_0$ \emph{и так} попадает в box. Тогда
клип ничего не меняет. Если выходит за границы --- клиппинг даёт
\emph{допустимое}, но \emph{не оптимальное} управление. В большинстве
рабочих режимов робота этого хватает.
\end{notebox}

\begin{codebox}[title={\filepath{mpc\_controller.py:288}}]
\begin{lstlisting}[language=Python]
def step(self, x, x_ref=None):
    if x_ref is None:
        dx = x
    else:
        dx = x - x_ref
    if self._solver == "qp":
        return self._solve_qp(dx)
    # Explicit clip
    u = -self.K_first @ dx
    return np.clip(u, self.u_min, self.u_max)
\end{lstlisting}
\end{codebox}

\subsection{\texttt{qp} (полный)}

При \texttt{control.mpc.solver: qp} в конфиге проект решает QP \emph{честно}:

\begin{codebox}[title={\filepath{mpc\_controller.py:303}}]
\begin{lstlisting}[language=Python]
def _solve_qp(self, dx):
    f = self.f_mat @ dx
    H = self.H
    u_lo = np.tile(self.u_min, self.N)
    u_hi = np.tile(self.u_max, self.N)
    bounds = list(zip(u_lo, u_hi))
    U0 = -np.linalg.solve(H, f)              # warm start
    U0 = np.clip(U0, u_lo, u_hi)
    res = minimize(lambda U: 0.5*U@H@U + f@U,
                   U0, jac=lambda U: H@U + f,
                   method="L-BFGS-B", bounds=bounds)
    return (res.x if res.success else U0)[:self.r]
\end{lstlisting}
\end{codebox}

Используется L-BFGS-B (scipy) с warm start от аналитического решения.
Делает то же, что OSQP, только медленнее. \emph{Включают только если
clip регулярно упирается в ограничения}.

\section{Выбор горизонта $N$}

\begin{ideabox}
\textbf{Большой $N$} (например, 50):
\begin{itemize}
  \item Регулятор «видит дальше» --- может заранее тормозить, чтобы
        не перелететь цель.
  \item QP становится больше: размер $\vect{U}$ растёт как $rN$,
        размер $\mat{H}$ --- как $(rN)^2$. Время на тик растёт нелинейно.
\end{itemize}

\textbf{Маленький $N$} (например, 1):
\begin{itemize}
  \item Решается быстро.
  \item Регулятор «близорукий», поведение ближе к ЛКР с малым весом
        конечного штрафа.
\end{itemize}

\textbf{Рабочий компромисс:} $N \cdot T_s \gtrsim 2 \tau_\text{slow}$, где
$\tau_\text{slow}$ --- самая медленная постоянная времени объекта. Для
нас $\tau_\text{slow} = \tau_v = 0.15$~с, $T_s = 0.02$~с $\Rightarrow$
$N \ge 15$. В дефолтах --- $N = 10$ или $20$ (см. \texttt{mps.horizon\_N}).
\end{ideabox}

\section{Терминальный штраф $\mat{P}_f$}

\begin{ideabox}
\textbf{Проблема конечного горизонта:} что мешает MPC сделать «хитрое»
движение, которое в окне $[0, N]$ выглядит выгодно, а после $N$
оставляет робот «в плохой позе» (далеко от цели)? Ничего, если стоимость
после $N$ мы игнорируем.

\textbf{Решение:} добавить терминальный штраф $\vect{x}_N\T \mat{P}_f\,
\vect{x}_N$, который аппроксимирует «стоимость до бесконечности» из
$\vect{x}_N$.

Лучшая такая аппроксимация --- решение DARE \eqref{eq:dare} для
$(\mat{A}_d, \mat{B}_d, \mat{Q}, \mat{R})$, то есть $\mat{P}_f = \mat{P}_\text{LQR}$.
Тогда MPC \emph{в пределе} $N \to \infty$ сходится к ЛКР, а на конечном
$N$ ведёт себя «как ЛКР, но с уважением к ограничениям».
\end{ideabox}

\begin{codebox}[title={\filepath{mpc\_controller.py:126}}]
\begin{lstlisting}[language=Python]
def _compute_terminal_penalty(self):
    """If Pf not given, solve DARE for a guaranteed-stable terminal cost."""
    return solve_discrete_are(self.Ad, self.Bd, self.Q, self.R)
\end{lstlisting}
\end{codebox}

\section{Слежение за reference}

До сих пор минимизировали $\vect{x}\T \mat{Q}\vect{x}$ --- штраф за
отклонение от \emph{нуля}. Но цель сценария --- $\vect{x}^\text{ref} \ne 0$
(например, $s_\text{ref}$ растёт линейно к $D$). Стандартный приём ---
работать в \emph{ошибке слежения} $\Delta\vect{x} = \vect{x} - \vect{x}^\text{ref}$:

\begin{equation}
  J = \sum_{i=0}^{N-1}\Big((\vect{x}_i - \vect{x}^\text{ref})\T \mat{Q}\,(\vect{x}_i - \vect{x}^\text{ref})
        + \vect{u}_i\T \mat{R}\vect{u}_i\Big)
        + (\vect{x}_N - \vect{x}^\text{ref})\T \mat{P}_f\,(\vect{x}_N - \vect{x}^\text{ref}).
  \label{eq:mpc_tracking}
\end{equation}

В коде --- именно так:

\begin{codebox}[title={\filepath{mpc\_controller.py:288}}]
\begin{lstlisting}[language=Python]
def step(self, x, x_ref=None):
    dx = x - x_ref  if x_ref is not None  else x
    ...
\end{lstlisting}
\end{codebox}

\section{Sanity-check замкнутой системы}

После \texttt{rebuild()} (Apply из UI) можно убедиться, что новый
регулятор устойчив:

\begin{codebox}[title={\filepath{mpc\_controller.py:332}}]
\begin{lstlisting}[language=Python]
def closed_loop_eigenvalues(self):
    """|λ| < 1 ⇔ closed-loop x_{k+1} = (Ad - Bd K_first) x is stable."""
    return np.linalg.eigvals(self.Ad - self.Bd @ self.K_first)

def is_stable(self):
    return bool(np.all(np.abs(self.closed_loop_eigenvalues()) < 1.0 - 1e-9))
\end{lstlisting}
\end{codebox}

Этот тест в UI делает «Validate matrices» --- кнопка перед Apply.

\section{Hot-reload между прогонами}

В МПС-учебном модуле пользователь \emph{редактирует} матрицы из UI и
жмёт Apply. На Pi прилетает \texttt{mps/matrices/set}, и нода вызывает
\texttt{mpc.rebuild()}: атомарно пересчитываются $\mat{\Phi}, \mat{\Gamma},
\mat{H}, \mat{K}_\text{first}$, при ошибке --- rollback к предыдущему
снимку. Это и есть «учебный пульт».

\begin{codebox}[title={\filepath{mpc\_controller.py:172}, фрагмент rebuild}]
\begin{lstlisting}[language=Python]
def rebuild(self, *, Ad=None, Bd=None, Q_diag=None, R_diag=None,
            N=None, u_min=None, u_max=None, Pf=None):
    snapshot = (self.Ad.copy(), ..., self.K_first.copy())
    try:
        # обновить поля, провалидировать формы, пересчитать matrices
        self._build_qp_matrices()
    except Exception:
        # Rollback on ANY failure
        (self.Ad, ..., self.K_first) = snapshot
        raise
\end{lstlisting}
\end{codebox}

\section{Что мы получили}

\begin{itemize}
  \item \textbf{MPC} --- решение QP-задачи на каждом тике с явными
        ограничениями.
  \item Через \emph{лифтинг} динамики QP сводится к канонической форме
        $\half\vect{U}\T \mat{H}\vect{U} + \vect{f}\T \vect{U}$.
  \item Без ограничений решение аналитическое, $\vect{u}_0 =
        -\mat{K}_\text{first}\vect{x}_0$. С ограничениями --- настоящая
        QP-задача, в проекте --- L-BFGS-B с warm start.
  \item Дефолтный режим --- \texttt{clip}: считаем $\mat{K}_\text{first}$,
        умножаем, клипуем. Быстро и в большинстве случаев оптимально.
  \item Терминальный штраф $\mat{P}_f$ --- решение DARE; делает короткий
        горизонт «бесконечно-устойчивым».
  \item Горизонт $N$ --- компромисс: $N T_s \gtrsim 2 \tau_\text{slow}$.
  \item \texttt{rebuild()} --- атомарное пересоздание контроллера с
        rollback на ошибку; используется при Apply из UI.
\end{itemize}

Теперь, наконец, всё готово, чтобы запустить настоящий сценарий.

% ============================================================
%  Глава 9 — Сценарий «D метров вперёд»
% ============================================================
\chapter{Сценарий «D метров вперёд»}
\label{ch:scenario}

\section{Что мы строим}

Соберём всё: модель ($\mat{A}_d, \mat{B}_d$), регулятор ($\mat{K}_\text{first}$
из MPC), измерение состояния ($s, v, \theta, \omega, e_\text{int}$),
оркестратор (FSM). Получится \textbf{сценарий} --- от нажатия кнопки в UI
до полного отчёта.

\begin{ideabox}
Сценарий --- это «контракт» между UI и роботом. Он говорит:
\emph{при каких параметрах робот считается успешно справившимся,
а при каких --- нет}. В МПС-модуле сценарий один: «проехать $D$~м
вперёд» (с опциональным предварительным разворотом, о нём ---
глава~\ref{ch:turn_phase}). Каждый прогон --- финитный, не цикл.
\end{ideabox}

\section{Постановка задачи}

\textbf{Дано:}
\begin{itemize}
  \item Дистанция $D \in (0, D_\text{max}]$~м, по умолчанию $2.0$.
  \item Целевая скорость $v_\text{target} \in (0, v_\text{max}]$~м/с,
        по умолчанию $0.15$.
  \item Целевой курс $\varphi \in [-\pi, \pi]$ (см. главу~\ref{ch:picker}).
        По умолчанию $0$ --- едем прямо.
\end{itemize}

\textbf{Требуется:} от старта (положение $s_\text{start}$, курс
$\theta_\text{start}$) переместиться в точку с относительной
координатой $s = D$, курс удерживать у $\varphi$.

\textbf{Допустимые исходы:}
\begin{center}
\begin{tabular}{lll}
\toprule
\textbf{Событие} & \textbf{Условие} & \texttt{status}\\
\midrule
Достижение цели & $s \ge D - 0.05$~м & \texttt{reached}\\
Тайм-аут        & $t > 3\,D / v_\text{target}$ & \texttt{timeout}\\
Принудит. abort & MQTT \texttt{mps/scenario/abort} & \texttt{aborted}\\
Расходимость    & $\|\vect{x}\| > $ bound, NaN/Inf & \texttt{error}\\
\bottomrule
\end{tabular}
\end{center}

\section{Reference}

MPC требует $\vect{x}^\text{ref}_k$ --- куда хотим попасть на каждом
шаге. Самая простая стратегия: \textbf{линейный разгон до $D$ за
номинальное время}.

\begin{formulabox}[title={Reference сценария}]
\begin{equation}
  s_\text{ref}(k) = \min\big(D,\;k\,\Delta t \cdot v_\text{target}\big),
  \qquad
  \vect{x}^\text{ref}_k = \begin{pmatrix} s_\text{ref}(k)\\ v_\text{target}\\ \varphi\\ 0\\ 0 \end{pmatrix}.
  \label{eq:reference}
\end{equation}
\end{formulabox}

Геометрия: $s_\text{ref}$ растёт линейно с угловым коэффициентом
$v_\text{target}$, пока не достигнет $D$ --- дальше «плато». Курс
$\theta_\text{ref} = \varphi$ постоянен. Угловая скорость
$\omega_\text{ref} = 0$ --- робот не должен крутиться, удерживая курс.

\begin{notebox}
Почему $v_\text{ref} = v_\text{target}$ \emph{константа}, а не «по
правилу»? Потому что в постоянной части $s_\text{ref}$ (после
достижения плато) скорость и так логически 0, но \textbf{штраф $\mat{Q}$
по $v$ должен быть маленький} --- иначе MPC будет «тормозить»
заранее, не доезжая до $D$. В \filepath{config.yaml} мы видим
$Q_\text{diag}[v] = 1$ против $Q_\text{diag}[s] = 10$ --- скорость
важна сильно меньше позиции.
\end{notebox}

\section{Полный тик регулятора}

Один полный шаг \texttt{mps\_node} (50~Гц):

\begin{enumerate}
  \item \textbf{Прочитать измерение.} Берём последний \texttt{x\_meas} из
        одометрии. Конвертируем в относительные координаты
        \eqref{eq:relative_state}.
  \item \textbf{Построить reference.} По текущему $k$ собираем
        $\vect{x}^\text{ref}_k$ из \eqref{eq:reference}.
  \item \textbf{Шаг MPC.} \texttt{u = mpc.step(x\_meas, x\_ref)}.
        Под капотом --- $\vect{u}_0 = -\mat{K}_\text{first}(\vect{x} -
        \vect{x}^\text{ref})$ + клиппинг.
  \item \textbf{Sanity-check.} Если в $\vect{u}$ или $\vect{x}$ есть
        NaN/Inf --- немедленный \texttt{status='error'}.
  \item \textbf{Жёсткий cap угловой скорости в forward-режиме.}
        $|u_\omega| \le 0.5$~рад/с (а не $2.0$ как глобально) --- защита
        от случайного «винта» при удержании курса.
  \item \textbf{Опубликовать.} \texttt{cmd\_vel} в MQTT
        (\texttt{linear\_x}, \texttt{angular\_z}); телеметрию \texttt{x, u, y}
        --- в UI.
  \item \textbf{Проверить условия завершения.}
        Достигли? Тайм-аут? Watchdog (нет одометрии 3 тика подряд)?
\end{enumerate}

\begin{codebox}[title={\filepath{pi\_nodes/nodes/mps\_node.py:436}}]
\begin{lstlisting}[language=Python]
def _tick_drive(self, run, x):
    phi = run.target_heading
    # Reference: ramp s_ref to D, hold v_target, hold heading φ.
    s_ref = min(run.distance, run.drive_t * run.v_target)
    x_ref = np.array([s_ref, run.v_target, phi, 0.0, 0.0])

    try:
        u = self._mpc.step(x, x_ref=x_ref)
    except Exception as exc:
        self._finish_run('error', f'mpc.step: {exc}')
        return

    if not (np.all(np.isfinite(u)) and np.all(np.isfinite(x))):
        self._finish_run('error', 'NaN/Inf in u or x')
        return

    # Hard omega cap in forward scenario.
    u[1] = max(-self._omega_max_fwd,
               min(self._omega_max_fwd, u[1]))

    self._publish_cmd_and_telemetry(run, x, u)
    # ... check reached, timeout, watchdog
\end{lstlisting}
\end{codebox}

\section{Lifecycle прогона}

\begin{center}
\begin{tikzpicture}[
  node distance=12mm and 18mm,
  state/.style={
    draw, rounded corners=3pt, thick, fill=blue!8,
    align=center, font=\small, minimum width=3cm, minimum height=10mm,
  },
  flow/.style={-{Stealth[length=2.5mm]}, thick},
]
  \node[state] (idle) {\textbf{IDLE}};
  \node[state, right=of idle] (run) {\textbf{DRIVE\_FORWARD\_MPS}\\
                                     tick @ 50 Гц};
  \node[state, right=of run, fill=green!8]  (fin) {\textbf{FINISHED}\\
                                                   cmd\_vel = 0 $\times 3$};
  \node[below=14mm of run, font=\footnotesize, align=center]
       (events) {$s \ge D - 0.05$ (\texttt{reached})\\
                 $t > 3D/v$ (\texttt{timeout})\\
                 \texttt{abort}\\
                 NaN, $\|x\| > $bound (\texttt{error})};
  \draw[flow] (idle) -- node[above, font=\scriptsize]{mps/scenario/run} (run);
  \draw[flow, dashed] (run) -- (events);
  \draw[flow] (run) -- node[above, font=\scriptsize]{одно из событий} (fin);
  \draw[flow] (fin.south) .. controls +(0,-2.5) and +(0,-2.5) .. (idle.south);
\end{tikzpicture}
\end{center}

Failsafe в FINISHED: моторам отправляется команда $[0, 0]$
\textbf{три раза} подряд. Это страховка на случай, если первый пакет
потеряется в MQTT --- мотор-нода имеет встроенный watchdog
\texttt{cmd\_vel\_timeout}, который сам стопнет двигатели, но дублирование
не помешает.

\section{Метрики качества}

После окончания прогона на основе массива телеметрии считаются
\textbf{метрики}. Они отвечают на вопрос «как хорошо отработал
регулятор». Без них «работает» не значит «хорошо».

\begin{center}
\begin{tabular}{lll}
\toprule
\textbf{Метрика} & \textbf{Формула} & \textbf{Что значит}\\
\midrule
\texttt{overshoot} & $\max_t s(t) - D$, обрезано $\ge 0$ & «Пролетели мимо»\\
\texttt{settling\_time} & первое $t$, после которого $|s-D| < 0.02$ м & «Сколько успокаивались»\\
\texttt{control\_energy} & $\sum_k \vect{u}_k\T \mat{R}\vect{u}_k \Delta t$ & «Расход управления»\\
\texttt{ss\_error} & $|s(t_\text{end}) - D|$ & «Установившаяся ошибка»\\
\texttt{peak\_v} & $\max_t |v(t)|$ & «Пиковая скорость»\\
\texttt{peak\_omega} & $\max_t |\omega(t)|$ & «Пиковая угловая скорость»\\
\bottomrule
\end{tabular}
\end{center}

\begin{ideabox}
\textbf{Учебный смысл метрик.} Можно специально подобрать «плохой»
$\mat{Q}$ или $\mat{R}$ так, чтобы:
\begin{itemize}
  \item получить большой \texttt{overshoot} (например, $r_v$ малая ---
        регулятор «топит педаль»);
  \item получить большой \texttt{settling\_time} ($r_v$ большая ---
        регулятор «медленный»);
  \item получить большую \texttt{ss\_error} ($q_s$ малая --- регулятору
        не критично, доехал ли).
\end{itemize}
Это и есть «учебный пульт» из спеки: пробуем разные веса, смотрим, что
получилось. Курсовая.
\end{ideabox}

\section{Safety}

Чтобы случайный «учебный пульт» не сломал железо, есть многоуровневая
защита:

\begin{enumerate}
  \item \textbf{Pre-validate}: запрос с $D > D_\text{max}$ или
        $v_\text{target} > v_\text{max}$ отбрасывается ещё на REST
        (HTTP 400) и ещё раз на Pi (\texttt{mps/error},
        \texttt{error\_type="precondition"}).
  \item \textbf{Watchdog по одометрии}: если 3 тика подряд нет
        \texttt{odom} --- прогон завершается ошибкой
        \texttt{watchdog}.
  \item \textbf{FSM-lock}: в \texttt{DRIVE\_FORWARD\_MPS} игнорируются
        voice/joystick/детекции мяча. Только \texttt{abort} переключает
        состояние.
  \item \textbf{Жёсткий cap $u_\omega$}: в forward-режиме $|u_\omega|
        \le 0.5$~рад/с (это меньше, чем общий $u_\text{max}[1] = 2$).
  \item \textbf{Failsafe}: при любом FINISHED тройной \texttt{cmd\_vel = 0}.
\end{enumerate}

\section{Симулятор}

Параллельно с роботом существует \emph{идеальный} симулятор ---
\filepath{compute\_node/mps\_runner.py}. Он гоняет ровно ту же
последовательность \texttt{mpc.step(x, x\_ref)}, но вместо одометрии
обновляет $\vect{x}$ по дискретной модели:

\[
  \vect{x}_{k+1} = \mat{A}_d\,\vect{x}_k + \mat{B}_d\,\vect{u}_k.
\]

Без шумов, без motor lag, без проскальзывания. Если запустить «Run on Sim»
из UI, прогон выполнится за $\sim 100$~мс прямо на бэкенде --- никакого
железа. Это нужно:

\begin{itemize}
  \item для UI --- быстрый итеративный подбор $\mat{Q}, \mat{R}$ без
        возни с роботом;
  \item для тестов --- юнит-тесты не требуют железа;
  \item для сравнения --- разница между sim и реальным прогоном
        показывает, насколько модель отстаёт от реальности.
\end{itemize}

\section{Что мы получили}

\begin{itemize}
  \item Сценарий определяет, \emph{что значит} «успешный прогон»:
        \texttt{reached}, \texttt{timeout}, \texttt{aborted}, \texttt{error}.
  \item Reference --- линейный разгон $s_\text{ref}$ до $D$ за
        номинальное время.
  \item Один тик --- 7 шагов от чтения одометрии до публикации
        \texttt{cmd\_vel}.
  \item Метрики --- объективная мера качества, считаются по концу
        прогона.
  \item Безопасность --- 5 уровней защиты от ломки железа учебным
        пультом.
  \item Симулятор --- идеальная модель, быстрая проверка без робота.
\end{itemize}

Это была \emph{одно}фазная картина. Но в текущей ветке мы добавили
\textbf{предварительный разворот} к выбранному курсу. Этим и
займёмся.

% ============================================================
%  Глава 10 — Двухфазный сценарий TURN → DRIVE
% ============================================================
\chapter{Двухфазный сценарий: TURN, затем DRIVE}
\label{ch:turn_phase}

\section{Вопрос}

В прошлой главе сценарий ехал прямо «вперёд от текущей ориентации».
А если пользователь \emph{выбрал} цель не прямо перед роботом, а
под углом $\varphi$ (например, $+45^\circ$ влево)? Что делать?

\begin{ideabox}
Можно было бы дать MPC сразу полную задачу: «доехать до точки $(D\cos\varphi,
D\sin\varphi)$, развернувшись по ходу». Но это плохая идея для линеаризованной
вокруг \emph{прямолинейного} движения модели: $\varphi = 45^\circ$ ---
далеко не «малое отклонение», линейная модель там врёт.

\textbf{Решение --- двухфазный сценарий:}
\begin{enumerate}
  \item \textbf{TURN.} Робот стоит на месте, поворачивается к курсу
        $\varphi$. После завершения --- курс точно $\varphi$, позиция
        не изменилась.
  \item \textbf{DRIVE.} Робот едет $D$ метров, удерживая курс $\varphi$.
        Это уже та же задача из главы~\ref{ch:scenario}, только
        $\theta_\text{ref} = \varphi$ вместо $0$.
\end{enumerate}
Между фазами --- мгновенный переход (без паузы), но \emph{с разными
ограничениями} на $u_\omega$ и \emph{другим} reference.
\end{ideabox}

\section{Фаза TURN}

\subsection{Reference}

В фазе TURN мы хотим, чтобы робот стоял ($s = 0$, $v = 0$) и развернулся
к курсу $\varphi$:

\begin{formulabox}[title={Reference фазы TURN}]
\begin{equation}
  \vect{x}^\text{ref}_\text{turn} = \begin{pmatrix} 0\\ 0\\ \varphi\\ 0\\ 0 \end{pmatrix}.
  \label{eq:turn_ref}
\end{equation}
\end{formulabox}

Подадим этот reference в \texttt{mpc.step(x, x\_ref)}. Полученное
$\vect{u}_\text{mpc}$ будет «оптимальным» --- но MPC не знает, что мы
хотим \emph{чистый разворот}, поэтому без дополнительных охранников
$u_v$ окажется не нулевым (особенно если $\varphi$ большой).

\subsection{Принудительные ограничения фазы}

В коде это решают \textbf{пост-обработкой}:

\begin{codebox}[title={\filepath{pi\_nodes/nodes/mps\_node.py:384}, фаза TURN}]
\begin{lstlisting}[language=Python]
def _tick_turn(self, run, x):
    phi = run.target_heading

    # Курс совпал с целью → переход в DRIVE.
    if abs(_normalize_angle(x[_THETA] - phi)) < self._turn_tol:
        run.phase = 'drive'
        run.drive_t = 0.0
        return True

    x_ref = np.array([0.0, 0.0, phi, 0.0, 0.0])
    u = self._mpc.step(x, x_ref=x_ref)

    # Чистое вращение: ход — в ноль; ω — свой (более высокий) кап.
    u[0] = 0.0
    u[1] = max(-self._omega_max_turn,
               min(self._omega_max_turn, u[1]))

    self._publish_cmd_and_telemetry(run, x, u)
    # ... watchdog, timeout
\end{lstlisting}
\end{codebox}

\textbf{Что происходит:}
\begin{enumerate}
  \item \texttt{u[0] = 0.0} --- принудительный ход в ноль. MPC мог
        захотеть двигаться, но мы не разрешаем.
  \item \texttt{u[1]} --- угловая скорость зажимается в
        $[-\omega_\text{max,turn}, +\omega_\text{max,turn}]$, где
        $\omega_\text{max,turn} = 1.0$~рад/с --- \emph{больше}, чем в
        DRIVE-фазе ($0.5$). На месте можно крутиться быстрее.
\end{enumerate}

\begin{notebox}
Почему \emph{разный} cap на $\omega$ в TURN и DRIVE? Потому что:
\begin{itemize}
  \item В TURN мы хотим быстро развернуться --- иначе зачем вообще
        фаза.
  \item В DRIVE \emph{большой} $\omega$ при движении --- это «винт»,
        робот съезжает с курса. Лучше держать узкий cap, чтобы
        регулятор был «нежным» в коррекциях.
\end{itemize}
Два параметра: \texttt{mps.scenario.omega\_max\_in\_turn = 1.0} и
\texttt{omega\_max\_in\_forward = 0.5}.
\end{notebox}

\subsection{Условие завершения TURN}

Переходим в DRIVE, когда курс совпал с целью с заданным допуском:

\begin{equation}
  \big|\,\text{wrap}_{[-\pi, \pi]}(\theta - \varphi)\,\big| < \varepsilon_\text{turn}.
\end{equation}

В коде $\varepsilon_\text{turn} = 0.05$~рад $\approx 2.9^\circ$ ---
поправка \texttt{mps.scenario.turn\_tolerance\_rad}.

\begin{warnbox}
\textbf{Wrap-around в $[-\pi, \pi]$ --- важно!} Если $\theta = +179^\circ$
и $\varphi = -179^\circ$, разность «по-наивному» $= 358^\circ$ --- огромная,
хотя физически они почти совпадают.

Поэтому используем функцию \texttt{\_normalize\_angle(a)}:
\begin{equation}
  \text{wrap}_{[-\pi, \pi]}(a) = ((a + \pi) \bmod 2\pi) - \pi.
\end{equation}
В коде:
\begin{lstlisting}[language=Python]
def _normalize_angle(a: float) -> float:
    return (a + math.pi) % (2 * math.pi) - math.pi
\end{lstlisting}

Без этого, в редких случаях, MPC получал бы reference-ошибку курса
$\Delta\theta \sim 2\pi$ --- умножал бы её на $\mat{K}$ --- получал бы
безумно большое $u_\omega$ --- упирался в насыщение и крутился в
обратную сторону.
\end{warnbox}

\subsection{Тайм-аут фазы TURN}

Чтобы при поломанной IMU TURN не зациклился навсегда, есть собственный
тайм-аут: $t_\text{turn} \le 10$~с (параметр
\texttt{mps.scenario.turn\_timeout\_s}). Не успели --- завершаем
прогон с \texttt{status='timeout'}.

\section{Переход TURN \texorpdfstring{$\to$}{->} DRIVE}

Когда условие завершения TURN сработало, \texttt{\_tick\_turn}
возвращает \texttt{True}, и в \texttt{\_tick} тем же тиком вызывается
\texttt{\_tick\_drive(run, x)}. Между ними:

\begin{itemize}
  \item не теряем тик (не пропускаем 20~мс на лоське);
  \item \texttt{run.phase = 'drive'} --- фаза переключилась;
  \item \texttt{run.drive\_t = 0.0} --- часы DRIVE-фазы стартуют с нуля,
        отдельно от общего $t$ (для ramp $s_\text{ref}$ и DRIVE-тайм-аута).
\end{itemize}

\begin{codebox}[title={\filepath{mps\_node.py:366}, общий tick}]
\begin{lstlisting}[language=Python]
def _tick(self):
    with self._lock:
        run = self._run
        if run is None or self._fsm_state != 'DRIVE_FORWARD_MPS':
            return
        x = self._x_meas.copy()
    # Относительные координаты.
    x[_S] = x[_S] - run.s_start
    x[_THETA] = _normalize_angle(x[_THETA] - run.theta_start)

    # Двухфазный сценарий.
    if run.phase == 'turn':
        if not self._tick_turn(run, x):
            return  # ещё крутимся / прогон завершён
        # TURN завершился ЭТИМ тиком → продолжаем в DRIVE
    self._tick_drive(run, x)
\end{lstlisting}
\end{codebox}

\section{Фаза DRIVE}

После переключения в DRIVE логика та же, что в главе~\ref{ch:scenario},
\emph{с одной поправкой}: курс reference --- $\varphi$, а не $0$.

\[
  \vect{x}^\text{ref}_\text{drive}(k) =
  \begin{pmatrix} s_\text{ref}(k)\\ v_\text{target}\\ \varphi\\ 0\\ 0 \end{pmatrix},
  \qquad
  s_\text{ref}(k) = \min(D, k\,\Delta t \cdot v_\text{target}).
\]

\begin{ideabox}
\textbf{Почему именно так, а не «следовать дуге»?} Робот в DRIVE
\emph{уже} ориентирован к $\varphi$ (фаза TURN об этом позаботилась).
Он едет вперёд в своей системе координат, а курс держит постоянным.
Это самая простая постановка, и она работает.

Если бы мы хотели ехать по \emph{дуге} (например, для плавного объезда
препятствия), то и reference, и линеаризация были бы сложнее: вместо
постоянной опорной $(v_0, 0)\T$ надо было бы линеаризовать вдоль дуги
с $\bar\omega \ne 0$. Это уже не учебный сценарий, а полноценная
trajectory tracking.
\end{ideabox}

\section{Timeline прогона}

\begin{center}
\begin{tikzpicture}[
  scale=1.0, font=\small,
  bar/.style={very thick, line cap=round},
]
  % Ось времени
  \draw[-{Stealth}] (-0.2, 0) -- (10.5, 0) node[right]{$t$, с};
  % Метки
  \foreach \x in {0,2,4,6,8,10}
    \draw (\x, -0.05) -- (\x, 0.05) node[above, font=\tiny]{\x};
  % TURN полоска
  \draw[bar, red!70!black]  (0, -0.6) -- (2.5, -0.6);
  \node[red!70!black] at (1.25, -0.95) {\textbf{TURN}};
  \node[font=\tiny, red!70!black] at (1.25, -1.25) {$\theta \to \varphi$};
  % DRIVE полоска
  \draw[bar, blue!70!black] (2.5, -0.6) -- (9.5, -0.6);
  \node[blue!70!black] at (6, -0.95) {\textbf{DRIVE}};
  \node[font=\tiny, blue!70!black] at (6, -1.25) {$s \to D$, $\theta = \varphi$};
  % Переход
  \draw[dashed, gray] (2.5, 0.3) -- (2.5, -1.5);
  \node[font=\scriptsize, gray] at (2.5, 0.55) {TURN$\to$DRIVE};
  % FINISHED
  \draw[bar, green!50!black] (9.5, -0.6) -- (10.0, -0.6);
  \node[green!50!black, font=\tiny] at (10.4, -0.95) {STOP};
\end{tikzpicture}
\end{center}

\textbf{Поля \texttt{\_RunState}, обслуживающие двухфазность:}

\begin{itemize}
  \item \texttt{run.t} --- монотонное \emph{суммарное} время от старта
        прогона. Растёт всегда. Используется для тайм-аута фазы TURN
        и как абсцисса телеметрии.
  \item \texttt{run.drive\_t} --- часы только DRIVE-фазы, стартуют с
        $0$ в момент переключения. Используются для ramp $s_\text{ref}$
        и DRIVE-тайм-аута.
  \item \texttt{run.phase} --- \texttt{'turn'} либо \texttt{'drive'}.
\end{itemize}

\section{Граничные случаи}

\begin{whatifbox}
\textbf{$\varphi = 0$.} Условие завершения TURN срабатывает на первом
же тике (если курс точно $0$): $\theta_\text{start} = 0 \Rightarrow x[\theta]
= 0 \Rightarrow |0 - 0| = 0 < \varepsilon_\text{turn}$. Сценарий
\emph{скипает} фазу TURN и идёт сразу в DRIVE. Это и есть «прямо вперёд»
из предыдущей главы.

\textbf{$\varphi = \pi$ (поворот назад).} Худший случай: робот должен
развернуться на $180^\circ$. Через wrap-around $\theta - \varphi$ всё
ещё корректно (нормализуется в $[-\pi, \pi]$), но при $\theta = 0$,
$\varphi = \pi$ ошибка ровно $\pm\pi$ --- знак $u_\omega$
неопределён. На практике первый тик MPC выберет какое-то направление,
робот начнёт крутиться, и дальше всё разрешится. Скорее всего проект
зафиксирует $\varphi \in (-\pi, \pi]$, не \emph{строгий} $\pi$, но валидация
в \texttt{\_on\_scenario\_run} допускает.

\textbf{Watchdog в TURN.} Если на каком-то тике нет одометрии $> 3$
тиков подряд, мы не можем достоверно судить, повернулся ли робот.
TURN завершается \texttt{error}, FSM возвращается в IDLE.
\end{whatifbox}

\section{Зачем это всё было нужно}

\begin{ideabox}
Думаем учебно: какие были \emph{альтернативы}?

\textbf{(А) Не разделять фазы, дать MPC сразу полную задачу с целевой
точкой $(D\cos\varphi, D\sin\varphi)$.} Тогда модель должна была бы
быть \emph{инвариантной} к ориентации --- то есть линеаризация вокруг
прямолинейного движения уже не работает; нужна нелинейная модель или
linearization gain scheduling.

\textbf{(Б) Сначала повернуться, потом ехать --- но через два разных
регулятора.} TURN --- простой П-регулятор на угле, DRIVE --- MPC.
Так у нас и было раньше (старый PID). Но это значит, что в коде
\emph{два} регулятора, и переключение между ними может давать всплеск
по \texttt{cmd\_vel}. Также --- два разных набора тюнинга, два моделирования.

\textbf{(В, выбранный) Один и тот же MPC, разные reference и капы
управления в каждой фазе.} Минусы: дополнительная сложность
переключения. Плюсы: один контроллер --- одна модель --- единое поведение.
Прозрачно для учебного эксперимента: те же $\mat{Q}, \mat{R}$ работают
в обеих фазах.
\end{ideabox}

\section{Что мы получили}

\begin{itemize}
  \item Сценарий стал \textbf{двухфазным}: TURN $\to$ DRIVE.
  \item Используется тот же MPC, но с разными reference и капами:
        TURN режим --- $u_v = 0$, $|u_\omega| \le 1.0$;
        DRIVE режим --- $u_v$ свободно, $|u_\omega| \le 0.5$.
  \item Переход --- мгновенный, тем же тиком, без пропуска
        управления.
  \item Угол $\varphi$ работает корректно через wrap-around в $[-\pi, \pi]$.
  \item Граничные случаи ($\varphi = 0$, watchdog) обрабатываются
        самой логикой.
\end{itemize}

Остался один вопрос: \textbf{откуда берётся $\varphi$}? Это
3D-пикер цели --- последняя интересная математика в этом учебнике.

% ============================================================
%  Глава 11 — 3D-пикер цели: математика угла
% ============================================================
\chapter{3D-пикер цели: математика}
\label{ch:picker}

\section{Вопрос}

Пользователь хочет указать курс $\varphi$, на который робот развернётся
перед движением. Как именно? Можно было бы поставить поле «введите угол
в градусах», но это \emph{неудобно} --- человек редко интуитивно знает,
как $+72^\circ$ соотносится с «вон туда, под небольшим углом влево».

\begin{ideabox}
\textbf{3D-пикер} --- решение из feat/mps-target-picker. На экране
показана 3D-сцена: робот в центре, окружность фиксированного радиуса
вокруг него, маркер на окружности. Пользователь \emph{кликает} в любую
точку «пола» сцены --- маркер переезжает на окружность в направлении
клика, а угол $\varphi$ автоматически вычисляется.

Это даёт мгновенную визуальную обратную связь и интуитивный ввод.
Под капотом --- немного линейной алгебры и одна функция \texttt{atan2}.
\end{ideabox}

\section{Геометрия и соглашения}

В 3D-сцене Three.js (React Three Fiber) система координат другая, чем
в мире робота. Аккуратно зафиксируем:

\begin{center}
\begin{tabular}{lll}
\toprule
& \textbf{Мир робота} & \textbf{Three.js}\\
\midrule
$x$ ось & вперёд от робота (curr. heading)  & вправо от камеры\\
$y$ ось & влево от робота & вверх (высота над полом)\\
$z$ ось & вверх & к камере (от сцены к зрителю)\\
\bottomrule
\end{tabular}
\end{center}

Робот «смотрит» вдоль своей $+x$ при yaw $= 0$. В Three угол сцены
сделан так, что пол --- плоскость $y = 0$ Three; камера слегка выше
смотрит вниз.

\textbf{Соответствие координат:} мир $(w_x, w_y)$ соответствует Three
$(w_x, h, -w_y)$, где $h$ --- высота над полом (для маркера ---
константа \texttt{MARKER\_HEIGHT}). Минус в третьей координате потому
что Three-$z$ направлен «к камере», а мировой $y$ --- «влево».

\begin{notebox}
Соглашение об углах: $\varphi$ --- относительный курс. $\varphi = 0$
--- прямо вперёд (вдоль текущего yaw робота); $\varphi > 0$ --- влево
(CCW против часовой); $\varphi < 0$ --- вправо. Это совпадает с
соглашением ROS и нашим robot frame в коде.
\end{notebox}

\section{Точка клика \texorpdfstring{$\to$}{->} угол}

\subsection{Формула}

В Three.js «raycast» от позиции мыши даёт точку клика на полу
$(x_3, z_3)$. Переводим в мировые $(w_x, w_y) = (x_3, -z_3)$. Угол ---
обычная $\operatorname{atan2}$:

\begin{formulabox}[title={Угол от точки клика}]
\begin{equation}
  \varphi = \operatorname{atan2}(w_y, w_x) = \operatorname{atan2}(-z_3, x_3).
  \label{eq:click_to_angle}
\end{equation}
\end{formulabox}

Функция \texttt{atan2} (есть в любой математической библиотеке)
аккуратно ставит знаки и выдаёт результат в $(-\pi, \pi]$. Заодно
не путается на $w_x = 0$ (где «наивный» $\arctan(w_y/w_x)$ дал бы
деление на ноль).

\begin{codebox}[title={\filepath{compute\_node/frontend/src/lib/targetAngle.ts}}]
\begin{lstlisting}[style=tsstyle]
// Точка клика на полу сцены (Three-координаты x, z) → относительный курс φ.
export function groundPointToAngle(
  threeX: number, threeZ: number
): number {
  // мир: wx = threeX, wy = -threeZ. φ = atan2(wy, wx).
  return Math.atan2(-threeZ, threeX)
}
\end{lstlisting}
\end{codebox}

\subsection{Почему именно \texttt{atan2}?}

\begin{ideabox}
\textbf{Альтернатива --- $\operatorname{atan}(w_y/w_x)$.} Беда:
\begin{itemize}
  \item Делит на ноль при $w_x = 0$ (точно «влево» или «вправо»).
  \item Не различает квадранты: $\arctan(1/1) = \arctan(-1/-1) = \pi/4$,
        хотя углы --- $45^\circ$ и $-135^\circ$.
\end{itemize}

\texttt{atan2} принимает оба аргумента \emph{отдельно}, использует знаки,
покрывает все 4 квадранта и не падает на нуле. Это стандартный приём в
графике и робототехнике.
\end{ideabox}

\subsection{Численные примеры}

\begin{examplebox}
$(w_x, w_y) = (1, 0)$: точка прямо перед роботом. $\varphi = \operatorname{atan2}(0, 1) = 0$ --- «прямо».

$(w_x, w_y) = (0, 1)$: точка слева от робота. $\varphi = \operatorname{atan2}(1, 0) = \pi/2 \approx 90^\circ$.

$(w_x, w_y) = (-1, 0)$: точка прямо за роботом. $\varphi = \operatorname{atan2}(0, -1) = \pi$ --- разворот на $180^\circ$.

$(w_x, w_y) = (-1, -1)$: точка справа-сзади. $\varphi = \operatorname{atan2}(-1, -1) = -3\pi/4 \approx -135^\circ$.
\end{examplebox}

\section{Угол \texorpdfstring{$\to$}{->} позиция маркера}

Обратное: имея $\varphi$ и фиксированный радиус $N$ (радиус окружности
в сцене), найти позицию маркера в Three-координатах.

\begin{formulabox}[title={Маркер на окружности радиуса $N$}]
\begin{equation}
  (x_3, h, z_3) = (N\cos\varphi,\; h_\text{marker},\; -N\sin\varphi).
  \label{eq:angle_to_pos}
\end{equation}
\end{formulabox}

Объяснение: в мировой системе точка на окружности радиуса $N$ под
углом $\varphi$ --- это $(N\cos\varphi, N\sin\varphi)$. Переходим в
Three: $z_3 = -w_y = -N\sin\varphi$, $x_3 = w_x = N\cos\varphi$.
Высота $h_\text{marker}$ --- константа над полом (чтобы маркер не
тонул в пол).

\begin{codebox}[title={\filepath{lib/targetAngle.ts}}]
\begin{lstlisting}[style=tsstyle]
export const MARKER_HEIGHT = 0.03

export function angleToMarkerPosition(
  angle: number,
  radius: number,
): [number, number, number] {
  return [radius * Math.cos(angle), MARKER_HEIGHT, -radius * Math.sin(angle)]
}
\end{lstlisting}
\end{codebox}

\subsection{Почему радиус \emph{фиксированный}?}

В предыдущей версии (одна фаза, без TURN) пикер позволял указать
\emph{любую точку} --- и расстояние, и направление. Это перегружало
интерфейс: пользователь думает «куда» И «насколько далеко» одновременно.

В двухфазной схеме \emph{дистанция $D$ задаётся отдельным полем},
а пикер отвечает только за \emph{направление}. Поэтому маркер всегда
ползёт по окружности фиксированного радиуса $N$ --- независимо от того,
куда конкретно кликнул пользователь. Это \emph{разделение
ответственности}: одна задача --- один контрол.

\section{Центральная мёртвая зона}

\begin{ideabox}
\textbf{Проблема:} если кликнуть \emph{рядом с центром}, маленькое
движение мыши даёт огромный скачок $\varphi$. У точки $(0.01, 0.01)$
угол $= \pi/4$, у точки $(0.01, -0.01)$ угол $= -\pi/4$ --- разница в
$90^\circ$ при микроперемещении в $1$~см.

\textbf{Решение:} в зоне радиуса $r_\text{dead} = 0.3\,N$ от центра
\emph{игнорируем} клик.
\end{ideabox}

\begin{codebox}[title={\filepath{lib/targetAngle.ts}}]
\begin{lstlisting}[style=tsstyle]
export const CENTER_DEADZONE_FRACTION = 0.3

export function groundPointRadius(
  threeX: number, threeZ: number
): number {
  return Math.hypot(threeX, threeZ)
}
\end{lstlisting}
\end{codebox}

В UI-логике: если \texttt{groundPointRadius < CENTER\_DEADZONE\_FRACTION * N},
клик \emph{не} обновляет $\varphi$.

\section{Подпись угла}

Финальный штрих --- человекочитаемая подпись маркера: «прямо», «+35°»,
«-40°»:

\begin{codebox}[title={\filepath{lib/targetAngle.ts}}]
\begin{lstlisting}[style=tsstyle]
export function formatHeadingLabel(angle: number): string {
  const deg = (angle * 180) / Math.PI
  if (Math.abs(deg) < 3) return 'прямо'
  const rounded = Math.round(deg)
  return rounded > 0 ? `+${rounded}°` : `${rounded}°`
}
\end{lstlisting}
\end{codebox}

Дед-зона по \emph{углу} ($\le 3^\circ$ показывается как «прямо») --- не
путать с дед-зоной по \emph{радиусу}. Это разные защиты:

\begin{itemize}
  \item Центральная (радиус) --- против скачков при микроперемещениях.
  \item Угловая (3°) --- косметика лейбла, чтобы не мигало
        $+1^\circ / -1^\circ$ при дрожании руки.
\end{itemize}

\section{Полный пайплайн пикера}

\begin{center}
\begin{tikzpicture}[
  node distance=8mm,
  pipestep/.style={
    draw, rounded corners=2pt, thick, fill=blue!8,
    align=center, font=\footnotesize, minimum width=4.5cm, minimum height=8mm,
  },
  flow/.style={-{Stealth[length=2mm]}, thick},
]
  \node[pipestep] (s1) {\textbf{1.} Клик мыши $(p_x, p_y)$};
  \node[pipestep, below=of s1] (s2) {\textbf{2.} Raycast $\to$ $(x_3, z_3)$ на полу};
  \node[pipestep, below=of s2] (s3) {\textbf{3.} Дед-зона: $\sqrt{x_3^2+z_3^2} < 0.3N$?};
  \node[pipestep, below=of s3] (s4) {\textbf{4.} $\varphi = \operatorname{atan2}(-z_3, x_3)$};
  \node[pipestep, below=of s4] (s5) {\textbf{5.} Маркер $(N\cos\varphi, h, -N\sin\varphi)$};
  \node[pipestep, below=of s5] (s6) {\textbf{6.} Лейбл «прямо» / «+35$^\circ$»};
  \node[pipestep, below=of s6, fill=green!8] (s7) {\textbf{7.} На «Run»: $\varphi$ в \texttt{target\_heading}};
  \draw[flow] (s1) -- (s2);
  \draw[flow] (s2) -- (s3);
  \draw[flow] (s3) -- node[right, font=\tiny]{нет: продолжаем} (s4);
  \draw[flow] (s4) -- (s5);
  \draw[flow] (s5) -- (s6);
  \draw[flow] (s6) -- (s7);
  \draw[flow] (s3.east) .. controls +(2,0) and +(2,-1) ..
              node[right, font=\tiny]{да: игнор} (s1.east);
\end{tikzpicture}
\end{center}

\section{От UI к роботу}

Когда пользователь жмёт «Run on Robot», значение $\varphi$ уходит в
REST-запрос:

\begin{codebox}[title={REST}]
\begin{verbatim}
POST /api/v1/mps/scenario/run
Content-Type: application/json

{
  "source": "robot",
  "request": {
    "distance": 2.0,
    "v_target": 0.15,
    "target_heading": 0.7854   // <-- phi (угол) из пикера, рад
  }
}
\end{verbatim}
\end{codebox}

\texttt{routers/mps.py} публикует это в MQTT \texttt{mps/scenario/run},
\texttt{mps\_node} читает, валидирует ($-\pi \le \varphi \le \pi$),
ставит \texttt{run.target\_heading = $\varphi$}, фаза \texttt{'turn'}
--- и поехали.

\section{Тесты математики пикера}

Файл \filepath{compute\_node/frontend/src/lib/targetAngle.test.ts}
(см. в проекте) покрывает:

\begin{itemize}
  \item Прямые тестовые случаи: $(1, 0) \to 0$, $(0, -1) \to \pi/2$ (помним
        Three's $z$ перевёрнут), $(-1, 0) \to \pi$ и т.д.
  \item Симметрию: \texttt{angleToMarkerPosition(groundPointToAngle(x, z) | radius = N)}
        возвращает точку на окружности радиуса $N$ \emph{в том же
        направлении}, что исходный клик.
  \item Граничные значения подписи: $\pm 2.9^\circ \to$ «прямо»,
        $\pm 3.1^\circ \to$ числовая подпись.
  \item Дед-зону по радиусу.
\end{itemize}

Юнит-тесты в jsdom без WebGL --- именно потому модуль \texttt{targetAngle.ts}
\emph{не зависит} от Three: вся «математика» вынесена в чистые функции,
а 3D-визуал --- отдельным компонентом \texttt{MpsTargetScene.tsx}.

\section{Что мы получили}

\begin{itemize}
  \item Пикер цели --- интуитивный визуальный ввод угла $\varphi$.
  \item Математика тривиальная: $\varphi = \operatorname{atan2}(-z_3, x_3)$
        для клика $\to$ угла, $(N\cos\varphi, h, -N\sin\varphi)$ для угла $\to$ позиции.
  \item Соглашения о координатах фиксированы и совпадают с тем, что
        потом получает Pi.
  \item Дед-зона по радиусу защищает от скачков, дед-зона по углу
        --- косметика.
  \item Юнит-тестируется без WebGL --- спасибо чистым функциям.
\end{itemize}

\medskip
\textbf{Это была вся «активная» математика модуля МПС.} Дальше осталась
ещё одна, инфраструктурная глава --- как матрицы из MATLAB попадают
в Python.

% ============================================================
%  Глава 12 — Реализация: MATLAB → YAML → Python
% ============================================================
\chapter{Реализация: путь чисел}
\label{ch:implementation}

\section{Вопрос}

Все предыдущие главы говорили: «матрица $\mat{A}$ имеет вот такой
вид», «$\mat{Q}$ выбирают так-то». Но как конкретное число $-6.667$ из
\eqref{eq:AB_numeric} попадает в Python-код, который шлёт команды
мотору? Где это происходит, и что меняется на каждом шаге?

\begin{ideabox}
В проекте --- классическая трёхэтапная схема:
\begin{enumerate}
  \item \textbf{Offline в MATLAB}: синтез матриц (модель, веса,
        проверка устойчивости).
  \item \textbf{Экспорт в YAML}: матрицы пишутся в \filepath{config.yaml}
        как обычные списки чисел.
  \item \textbf{Online в Python}: \texttt{mps\_node} читает YAML при старте,
        строит \texttt{StateSpaceModel} и \texttt{MPCController}, гонит
        тики.
\end{enumerate}
Плюс \textbf{live-канал} от UI: пользователь меняет матрицы в браузере
$\to$ MQTT \texttt{mps/matrices/set} $\to$ \texttt{mpc.rebuild()} на лету.
\end{ideabox}

\section{Этап 1: MATLAB (\filepath{matlab/main.m})}

В \filepath{matlab/} лежат оффлайн-скрипты для синтеза матриц.
Логика стандартная:

\begin{enumerate}
  \item Загрузить параметры объекта ($v_0, \tau_v, \tau_\omega$) из конфига.
  \item Построить $\mat{A}, \mat{B}$ по формулам \eqref{eq:A_continuous},
        \eqref{eq:B_continuous}.
  \item Проверить управляемость: \texttt{rank(ctrb(A, B)) == n}.
  \item Выбрать $\mat{Q}, \mat{R}$.
  \item Решить DARE (MATLAB-функция \texttt{dlqr}/\texttt{dare}) ---
        получить $\mat{P}, \mat{K}$.
  \item Визуализировать step response: убедиться, что замкнутая
        система устойчива и качественно ведёт себя.
  \item Экспортировать всё в YAML.
\end{enumerate}

\begin{codebox}[title={\filepath{matlab/main.m} (упрощённый эскиз)}]
\begin{lstlisting}[style=matlabstyle]
% Параметры объекта
v0    = 0.20;
tau_v = 0.15;
tau_w = 0.10;
Ts    = 0.02;     % период МПС-цикла, 50 Гц

% Непрерывные матрицы
A = [0 0 0    1         0;
     0 0 v0   0         0;
     0 0 0    0         1;
     0 0 0   -1/tau_v   0;
     0 0 0    0        -1/tau_w];
B = [0 0; 0 0; 0 0; 1/tau_v 0; 0 1/tau_w];

% Управляемость
assert(rank(ctrb(A, B)) == 5);

% Дискретизация
sysd  = c2d(ss(A, B, eye(5), zeros(5,2)), Ts, 'zoh');
Ad = sysd.A;  Bd = sysd.B;

% Веса
Q = diag([10 10 5 1 1]);
R = diag([1 1]);

% LQR / DARE
[K, P, ~] = dlqr(Ad, Bd, Q, R);

% Проверка устойчивости замкнутой
assert(all(abs(eig(Ad - Bd*K)) < 1));

% Экспорт
export_to_yaml(A, B, Q, R, K, P, 'config.yaml');
\end{lstlisting}
\end{codebox}

\begin{notebox}
MATLAB используется потому что: (1) исторически он стандарт ТАУ-курсов,
(2) \texttt{c2d}, \texttt{dlqr}, \texttt{step}, графики --- встроенные
и предсказуемые, (3) учебник Козлова пишет в MATLAB-нотации.

Альтернатива --- Python+\texttt{scipy.signal}+\texttt{matplotlib}. Сейчас
проект использует MATLAB для синтеза и Python --- для онлайн-управления.
В \filepath{tools/} есть Python-эквивалент для тех, у кого нет лицензии
MATLAB.
\end{notebox}

\section{Этап 2: YAML --- источник правды}

\filepath{config.yaml} --- общая «истина» для Python и UI. В нём
\emph{две независимые} секции для двух моделей:

\begin{codebox}[title={\filepath{config.yaml}, секция \texttt{control} (легаси)}]
\begin{verbatim}
control:
  plant:
    Ts:     0.05
    v0:     0.20
    tau_v:  0.15
    tau_w:  0.10
  matrices:
    A: [[0, 0, 0,    1,         0],
        [0, 0, 0.2,  0,         0],
        [0, 0, 0,    0,         1],
        [0, 0, 0,   -6.667,     0],
        [0, 0, 0,    0,        -10]]
    B: [[0,      0],
        [0,      0],
        [0,      0],
        [6.667,  0],
        [0,     10]]
    K_mpc: [...]    # предвычисленное усиление
    Pf:    [...]    # терминальный штраф
  weights:
    Q_diag: [10.0, 10.0, 5.0, 1.0, 1.0]
    R_diag: [1.0, 1.0]
  mpc:
    horizon_N: 10
    solver:    "clip"
  limits:
    u_min: [-0.30, -2.0]
    u_max: [+0.30, +2.0]
\end{verbatim}
\end{codebox}

\begin{codebox}[title={\filepath{config.yaml}, секция \texttt{mps} (учебная)}]
\begin{verbatim}
mps:
  plant:
    Ts: 0.02       # 50 Гц
  matrices:
    A: [...]      # непрерывная (контракт api.md)
    B: [...]      # непрерывная
    C: [...]      # для UI визуализации (по умолч. I)
    D: [...]      # для UI визуализации (по умолч. 0)
  weights:
    Q_diag: [...]
    R_diag: [...]
  horizon_N: 20
  limits:
    u_min: [...]
    u_max: [...]
  scenario:
    distance_max:           5.0
    v_target_max:           0.30
    omega_max_in_forward:   0.5
    omega_max_in_turn:      1.0
    turn_tolerance_rad:     0.05
    turn_timeout_s:         10.0
    default_distance:       2.0
    default_v_target:       0.15
\end{verbatim}
\end{codebox}

\begin{warnbox}
\textbf{Две секции не пересекаются.} \texttt{control.*} --- для старого
PID-/легаси-MPC контура; \texttt{mps.*} --- для учебного МПС-модуля.
\emph{Не путать}: легаси-секция содержит \emph{дискретные} матрицы и
готовое усиление \texttt{K\_mpc}, а МПС-секция --- \emph{непрерывные}
матрицы, и ZOH-дискретизация делается в Python при загрузке.

Если случайно подать \texttt{mps}-матрицы в легаси-контур (или наоборот),
размерности совпадут (по $5\times 5$), но смысл компонент состояния --- разный,
и MPC «закрутит» робота не туда. Так и было в коммите \texttt{8521742}
до его починки.
\end{warnbox}

\section{Этап 3: Python (\texttt{StateSpaceModel}, \texttt{MPCController})}

При старте \texttt{mps\_node} читает \texttt{config.yaml} и собирает
plant + MPC:

\begin{codebox}[title={\filepath{pi\_nodes/nodes/mps\_node.py:180}, метод \texttt{\_build\_from\_config}}]
\begin{lstlisting}[language=Python]
def _build_from_config(self):
    A = self._cfg('mps.matrices.A', None)
    B = self._cfg('mps.matrices.B', None)
    C = self._cfg('mps.matrices.C', None)
    D = self._cfg('mps.matrices.D', None)
    Q_diag = self._cfg('mps.weights.Q_diag', None)
    R_diag = self._cfg('mps.weights.R_diag', None)
    N = self._cfg('mps.horizon_N', None)
    u_min = self._cfg('mps.limits.u_min', None)
    u_max = self._cfg('mps.limits.u_max', None)
    if any(v is None for v in (A, B, Q_diag, R_diag, N, u_min, u_max)):
        raise RuntimeError('config.yaml: секция mps: неполна')

    A = np.asarray(A, dtype=float)
    B = np.asarray(B, dtype=float)
    Ad, Bd = zoh_discretize(A, B, self._mps_ts)
    plant = StateSpaceModel(Ad=Ad, Bd=Bd, Cd=C, Dd=D, Ts=self._mps_ts)
    mpc = MPCController(
        Ad=Ad, Bd=Bd,
        Q=np.diag(np.asarray(Q_diag, dtype=float)),
        R=np.diag(np.asarray(R_diag, dtype=float)),
        N=int(N),
        u_min=np.asarray(u_min, dtype=float),
        u_max=np.asarray(u_max, dtype=float),
        solver='clip',
    )
    return plant, mpc
\end{lstlisting}
\end{codebox}

Что важно отметить:

\begin{itemize}
  \item Матрицы из YAML --- \emph{непрерывные}; ZOH-дискретизация
        происходит \emph{здесь}, в одной точке кода.
  \item \texttt{Ad, Bd} затем используются \emph{и} в plant
        (\texttt{StateSpaceModel}), \emph{и} в MPC --- единый источник
        правды. Симулятор на ноуте делает то же самое (\texttt{compute\_node/
        mps\_runner.py: \_build\_controller}) --- три места, одна логика.
  \item \texttt{solver='clip'} зашит --- учебный, быстрый, не требует
        scipy на Pi.
\end{itemize}

\section{Live Apply: четвёртый канал}

Сверх старта-из-конфига есть \textbf{live-канал}: пользователь
редактирует $\mat{A}, \mat{B}, \mat{Q}, \mat{R}$ в UI и жмёт Apply.

\subsection{Поток данных}

\begin{center}
\begin{tikzpicture}[
  node distance=8mm,
  pipestep/.style={
    draw, rounded corners=2pt, fill=blue!8,
    align=center, font=\footnotesize, minimum width=4.5cm, minimum height=8mm,
  },
  flow/.style={-{Stealth[length=2mm]}, thick},
]
  \node[pipestep] (s1) {UI: \texttt{MatrixEditor} (React)};
  \node[pipestep, below=of s1] (s2) {POST \texttt{/api/v1/mps/matrices/apply}};
  \node[pipestep, below=of s2] (s3) {\filepath{routers/mps.py}\\publish MQTT};
  \node[pipestep, below=of s3] (s4) {Pi: \texttt{mps\_node.\_on\_matrices\_set}};
  \node[pipestep, below=of s4] (s5) {\texttt{plant.reload()} + \texttt{mpc.rebuild()}};
  \node[pipestep, below=of s5, fill=green!8] (s6) {publish \texttt{mps/matrices/applied}};
  \node[pipestep, below=of s6] (s7) {UI: \texttt{state.mps.applied = matrices}};
  \foreach \a/\b in {s1/s2, s2/s3, s3/s4, s4/s5, s5/s6, s6/s7}
    \draw[flow] (\a) -- (\b);
\end{tikzpicture}
\end{center}

Ключевые моменты:

\begin{itemize}
  \item Запрос отбрасывается, если прогон идёт (\texttt{is\_running ==
        True}). Apply --- \emph{между} прогонами.
  \item На Pi A, B приходят в непрерывной форме --- ZOH-дискретизируем
        тем же \texttt{zoh\_discretize}, что и на старте.
  \item \texttt{rebuild()} --- атомарный: snapshot всех полей, при
        ошибке валидации (например, $\mat{R}$ не положительная) ---
        rollback к снимку. См. главу~\ref{ch:mpc}, §rebuild.
  \item Ack \texttt{mps/matrices/applied} приходит обратно в UI ---
        пользователь видит «зелёную галочку».
\end{itemize}

\subsection{Что валидирует Pi}

\begin{enumerate}
  \item Формы матриц: $\mat{A}$ квадратная, $\mat{B}$ совместима с
        $\mat{A}$, $\mat{Q}$ диагональная нужного размера, и т.д.
  \item Положительность: $r_i > 0$. Иначе DARE не сойдётся.
  \item Размер $N \ge 1$, $u_\text{min} < u_\text{max}$.
  \item Замкнутая система устойчива (опционально, в \texttt{mpc.is\_stable()}).
\end{enumerate}

При нарушении публикуется \texttt{mps/error}, тип
\texttt{error\_type="precondition"}. UI показывает алерт.

\section{Согласование sim и робота}

\begin{ideabox}
\textbf{Кошмарный сценарий:} симулятор и робот строят MPC \emph{по-разному}.
В sim матрицы дискретизируются с $T_s = 0.02$, на роботе --- $0.05$.
Или симулятор берёт $Q$ из YAML, а робот --- свои дефолты. Тогда
\emph{прогон в sim} перестаёт быть предсказанием реального прогона ---
учебный модуль теряет смысл.

\textbf{Решение:} \emph{три} места, где строится контроллер
(на старте Pi, при Apply на Pi, в симуляторе), должны быть \textbf{строго
идентичны} по логике. Это инвариант, который тесты охраняют
(\filepath{tests/test\_mps\_runner.py}, \filepath{tests/test\_mps\_node.py}).
\end{ideabox}

\section{Тесты}

\begin{center}
\begin{tabular}{lll}
\toprule
\textbf{Файл} & \textbf{Что проверяет} & \textbf{Где живёт}\\
\midrule
\texttt{test\_mps\_runner.py}     & Sim: propagation, reach, метрики & \filepath{tests/}\\
\texttt{test\_mps\_router.py}     & REST/WS контракты, history       & \filepath{tests/}\\
\texttt{test\_mps\_node.py}       & MQTT handlers, FSM, abort        & \filepath{tests/}\\
\texttt{test\_state\_space\_extended.py} & Cd, Dd, reload(), output() & \filepath{tests/}\\
\texttt{test\_mpc\_controller.py} & rebuild(), горизонт, веса, rollback & \filepath{tests/}\\
\texttt{targetAngle.test.ts}      & Математика пикера (jsdom)        & \filepath{compute\_node/frontend/}\\
\texttt{MpsPage.test.tsx}         & UI flow: пикер, Run, история     & \filepath{compute\_node/frontend/}\\
\bottomrule
\end{tabular}
\end{center}

Coverage таргет --- 85\% по проекту. CI на каждый PR в feat/mps
прогоняет все эти суиты, плюс \emph{smoke-тест на реальном Pi} (A10)
--- руками, по чек-листу.

\section{Где какие риски}

\begin{notebox}
Учебная природа модуля = вольный экспериментатор у руля. Места, где
\emph{реально} может «не пойти», и куда смотреть в первую очередь:

\begin{itemize}
  \item \textbf{$\mat{R}$ с нулём.} DARE не сойдётся, \texttt{rebuild()}
        упадёт. Pi ответит \texttt{precondition}.
  \item \textbf{$\mat{Q}$ с большими разбросами весов.} Может дать
        численно плохо обусловленную $\mat{H}$, $\mat{K}$ будет
        огромным $\to$ робот будет дёргаться. Pre-validate этого
        не ловит, ловит \texttt{is\_stable()}.
  \item \textbf{Несовпадение секций конфига и трактовки.} См. warning
        выше про коммит \texttt{8521742}.
  \item \textbf{Watchdog по одометрии.} Если выдернуть USB у Pi от
        мотор-контроллера во время прогона --- \texttt{watchdog}
        отработает за $3 \cdot T_s = 60$~мс.
\end{itemize}
\end{notebox}

\section{Что мы получили}

\begin{itemize}
  \item Числа модели проходят через три точки: MATLAB (синтез) ---
        YAML (источник правды) --- Python (онлайн).
  \item Кроме старта-из-конфига, есть live-канал: UI \texttt{->} MQTT
        \texttt{->} \texttt{mpc.rebuild()}. Между прогонами, атомарно,
        с rollback.
  \item Три места построения контроллера (Pi-start, Pi-Apply, simulator)
        обязаны быть идентичны --- это инвариант.
  \item Тестов --- семь файлов, покрывающих все три уровня (юнит на
        backend, юнит на frontend, smoke на Pi).
\end{itemize}

\medskip
\textbf{Поздравляем.} От «нажал кнопку» до ШИМ-сигнала в моторе ---
вся математика, вся архитектура, все подводные камни. Если ты
читал внимательно, ты теперь можешь:

\begin{enumerate}
  \item Объяснить другому, что значит «вектор состояния».
  \item Вывести $\mat{A}, \mat{B}$ из физической постановки.
  \item Понимать, что делает каждая строка \texttt{mps\_node.\_tick()}.
  \item Предсказать, как изменится поведение робота при изменении
        $\mat{Q}, \mat{R}, N$.
  \item Найти и починить баг, аналогичный коммиту \texttt{8521742}.
\end{enumerate}

\textbf{Удачи с курсовой.}


\end{document}
